\documentclass{article}

\usepackage[margin=1.6in]{geometry}

\title{Summary of essential results in the theory of topological tensor
products and nuclear spaces}
\author{A. Grothendieck}
\date{1952}

\usepackage{amssymb,amsmath,amscd}
\usepackage{hyperref}
\usepackage{xcolor}
\hypersetup{colorlinks,linkcolor={blue!50!black},citecolor={blue!50!black},urlcolor={blue!80!black}}
\usepackage{enumerate}
\usepackage{tikz-cd}
\usepackage{booktabs}
\usepackage{mathtools}

% for pandoc tables
\usepackage{calc,array}
\usepackage{longtable}
%

\usepackage{mathrsfs}
\usepackage{fouriernc}

\providecommand{\tightlist}{%
  \setlength{\itemsep}{0pt}\setlength{\parskip}{0pt}}


%% Header and footer %%

\usepackage{fancyhdr}
\usepackage{lastpage}
\usepackage{xstring}
\pagestyle{fancy}
\fancypagestyle{plain}{}
\fancyhf{}
\lhead{\footnotesize\nouppercase\leftmark}
\cfoot{\small\thepage\ of \pageref*{LastPage}}
% Git commit hash for server builds
\newif\ifserver
\serverfalse
\lfoot{\footnotesize\ifserver{Git commit: \href{https://github.com/thosgood/translations/commit/GitCommitHashVariable}{GitCommitHashVariable}}\fi}


%% Theorem environments %%

\usepackage{amsthm}

\newenvironment{itenv}[1]
  {\phantomsection\par\smallskip\noindent\textbf{#1.}\itshape}
  {\par\smallskip}

\newenvironment{rmenv}[1]
  {\phantomsection\par\smallskip\noindent\textbf{#1.}\rmfamily}
  {\par\smallskip}

\newenvironment{titenv}[1]
  {\phantomsection\par\noindent\textbf{#1.}\itshape}
  {\par}

\newenvironment{trmenv}[1]
  {\phantomsection\par\noindent\textbf{#1.}\rmfamily}
  {\par}


%% Shortcuts %%

\newcommand{\oldpage}[1]{\marginpar{\footnotesize$\Big\vert$ \textit{p.~#1}}}


%% Bibliography %%

\usepackage[backend=biber,maxnames=99,maxalphanames=4]{biblatex}
\addbibresource{AIF-4-1952-73.bib}


%% Document %%

\usepackage{embedall}
\begin{document}

\maketitle
\thispagestyle{fancy}

\renewcommand{\abstractname}{Translator's note}

\setcounter{footnote}{0}

% N.B. the TOC is now found in _translator-note.Rmd
% \tableofcontents


%% Content %%

\providecommand{\scr}[1]{{\mathscr{#1}}}
\renewcommand{\cal}[1]{{\mathcal{#1}}}
\renewcommand{\frak}[1]{{\mathfrak{#1}}}
\renewcommand{\geq}{\geqslant}
\renewcommand{\leq}{\leqslant}

\providecommand{\BB}{\mathrm{B}}
\providecommand{\LL}{\mathrm{L}}
\providecommand{\JJ}{\mathrm{J}}
\providecommand{\CC}{\mathrm{C}}
\providecommand{\RR}{\mathbb{R}}
\renewcommand{\ll}{\ell}
\providecommand{\sBB}{\mathscr{B}}
\providecommand{\sDD}{\mathscr{D}}
\providecommand{\sEE}{\mathscr{E}}
\providecommand{\sHH}{\mathscr{H}}
\providecommand{\sLL}{\mathscr{L}}
\providecommand{\sOO}{\mathscr{O}}
\providecommand{\sSS}{\mathscr{S}}
\providecommand{\DD}{(\mathscr{D})}
\providecommand{\DDp}{(\mathscr{D}')}
\providecommand{\DF}{(\mathscr{DF})}
\providecommand{\FF}{(\mathscr{F})}
\providecommand{\LF}{(\mathscr{LF})}
\providecommand{\MM}{(\mathscr{M})}
\providecommand{\hotimes}{\hat{\otimes}}
\providecommand{\scr}{\mathscr}
\providecommand{\dd}{\mathrm{d}}
\providecommand{\lotimes}{\bar{\otimes}}
\providecommand{\potimes}[1][p]{\overset{(#1)}{\otimes}}
\providecommand{\spotimes}[1][p]{\overset{[#1]}{\otimes}}
\providecommand{\Tr}{\operatorname{Tr}}
\providecommand{\hhotimes}{\hat{\hat{\otimes}}}

\tableofcontents

\section*{Introduction}\label{introduction}
\addcontentsline{toc}{section}{Introduction}

\subsection*{Subject}\label{subject}
\addcontentsline{toc}{subsection}{Subject}

\oldpage{73}

This article aims to give a summary, without proofs, of the principal
results found in my work ``Produits tensoriels topologiques et espaces
nucléaires'', which will be published in the \emph{Memoirs of the Amer.
Math. Society} \autocite{PTT} (and which I will refer to as {[}PTT{]}).
The main concern throughout {[}PTT{]} was that of being exhaustive, both
in terms of studying all the questions raised by the topics covered, as
well as trying to state the more difficult results as theorems that were
as general as possible. This work was also very dense, and the important
simple ideas risked being sometimes hidden behind technical details.
This is why this bowdlerised summary is possibly useful in giving a more
assimilable outline of the theory. Some extra comments, interesting but
not necessary for the general understanding of this summary, as well as
some hints for certain proofs, have been placed between stars, \(\star\)
like this \(\star\).

The importance of topological tensor products shows itself in many
different settings:

\oldpage{74}

\begin{enumerate}
\def\labelenumi{\alph{enumi}.}
\item
  The notion of the topological tensor product forms the foundations of
  a simple and general formulation of \emph{Fredholm theory}, including,
  alongside the classical case of an integral operator defined by a
  continuous kernel, many other operators that are defined in the most
  important functional spaces.\footnote{Such a formulation of Fredholm
    theory seems to have appeared for the first time in A. Ruston,
    ``Direct product of Banach spaces and linear functional equations'',
    \emph{Proc. of the London Math. Soc.} \textbf{3} (1951), 1. My work
    on this subject was conceived independently of his (in the autumn of
    1951), and is rather different.}
\item
  The many variants of the notion of topological tensor product give
  rise, by duality, to the definition of many remarkable classes of
  bilinear forms and linear operators, whose study is only just barely
  covered in {[}PTT, chap.~I, §4{]}. In particular, the techniques
  introduced there, conveniently systematised and exploited, allow us to
  obtain entirely unexpected results in the \emph{theory of linear
  transformations between the spaces \(L^1\), \(L^2\), and
  \(L^\infty\)}, and their topological-vectorial analogues (these
  results being, as of yet, not definitive, and thus unpublished). I
  might return to this subject, and restrict myself to explaining, in a
  rather different way, the systematic work of von Neumann--Schatten on
  the remarkable classes of compact operators in a Hilbert space
  \autocite[chap.~4]{8}.
\item
  From the point of view of this current work, the most important
  application of topological tensor products is the theory of
  \emph{nuclear spaces}. We explain this theory, generalise it, and make
  precise the famous ``theory of kernels'' of L. Schwartz, and further
  discover new properties, even for the most classical of spaces. Here,
  the topological tensor calculus is the most simple, since the majority
  of variants of the notion of topological tensor product coincide, and
  their properties thus sum. For now, there are not many applications of
  the general theorems that we obtain to specific theories. The most
  interesting seems to be a topological-vectorial variant of the
  ``Künneth theorem'', giving the homology of a complex defined as the
  tensor product of two complexes, a variant which seems useful in
  topological algebra.
\item
  Generally, it seems to me that the notions of topological tensor
  product are perfect for giving a suggestive and manageable
  \emph{language} that would be good to use in many situations in
  functional analysis, especially since we have theorems (some of which
  are non-trivial) at our disposition from which we can benefit. I hope
  that this summary (or, better, {[}PTT{]}) will succeed in giving the
  reader a similar impression, before the publication of the articles
  promised above.
\end{enumerate}

\subsection*{Terminology and notation}\label{terminology-and-notation}
\addcontentsline{toc}{subsection}{Terminology and notation}

Generally we follow the terminology and notation of \autocite{3}, apart
from the fact that we call the semi-reflexive spaces of \autocite{3}
\emph{reflexive}. We only consider, unless otherwise stated, spaces that
are \emph{locally convex} and \emph{separated}; by ``quotient space'' of
a space \(E\), we mean the quotient of \(E\) by a \emph{closed} vector
subspace. The \emph{dual} of \(E\), written \(E'\), is assumed to be,
unless otherwise stated, endowed with the strong topology (i.e.~the
topology of bounded convergence). \oldpage{75} The dual of \(E'\), or
the \emph{bidual} of \(E\), written \(E''\), is assumed to be, unless
otherwise stated, endowed with the topology given by uniform convergence
on the equicontinuous subsets of \(E'\), which induces the original
topology on \(E\). We will eventually need to appeal to certain notions
defined and studied in \autocite{6}, most notably that of a
\emph{\((\mathscr{DF})\)-space}. For our purposes here, it will suffice
to know that the dual of a \((\mathscr{F})\)-space is a
\((\mathscr{DF})\)-space; that every normed space is a
\((\mathscr{DF})\)-space; and that the dual of a
\((\mathscr{DF})\)-space is an \((\mathscr{F})\)-space.

Let \(E\), \(F\), and \(G\) be locally convex spaces. Denote by
\(\mathrm{B}(E,F;G)\) (resp. \(\mathscr{B}(E,F;G)\)) the space of
continuous bilinear maps (resp. of separately continuous bilinear maps,
i.e.~linear with respect to each variable) from \(E\times F\) to \(G\).
Denote by \(\mathrm{L}(E;F)\) the space of continuous linear maps from
\(E\) to \(F\). Denote by \(\mathscr{B}_e(E'_s,F'_s)\) the space of
separately continuous bilinear forms on the product of the weak duals
\(E'_s\) and \(F'_s\) of \(E\) and \(F\), endowed with the
\emph{biequicontinuous topology}, i.e.~the topology given by uniform
convergence on the products of an equicontinuous subset of \(E'\) with
equicontinuous subset of \(F'\). This space is complete if and only if
the spaces \(E\) and \(F\) are complete.

We define a \emph{bounded} (resp. \emph{compact}, resp. \emph{weakly
compact}) \emph{linear map} from \(E\) to \(F\) to be a linear map from
\(E\) to \(F\) that sends a suitable neighbourhood of \(0\) to a bounded
(resp. relatively compact, resp. relatively weakly compact) subset of
\(F\).

For short\footnote{This terminology was suggested to me by R.E. Edwards.
  \emph{{[}Trans.{]} The seemingly more popular terminology these days
  is to say ``absolutely convex'' instead of ``disked'', and to speak of
  the ``absolutely convex hull'' instead of the ``disked hull''.}}, if
\(E\) is a vector space, we say \emph{disk} or \emph{disked set in
\(E\)} to mean a convex and circled (a.k.a. balanced) subset of \(E\).
If \(E\) is a locally convex space, and \(A\) a bounded disk in \(E\),
then we denote by \(E_A\) the vector space generated by \(A\), and
endowed with the norm \(\|x\|_A=\inf_{x\in\lambda A}|\lambda|\). If
\(A\) is closed, then the unit ball of \(E_A\) is \(A\). If \(A\) is
complete, then \(E_A\) is complete. If \(V\) is a disked neighbourhood
of \(0\) in \(E\), then \(E_V\) denotes the normed space given by
passing to the quotient under the semi-norm
\(\|x\|_V=\inf_{x\in\lambda V}|\lambda|\).

Recall that a locally convex space is said to be \emph{quasi-complete}
if its closed bounded subsets are complete, \emph{barrelled} (resp.
\emph{quasi-barrelled}) if the bounded subsets of its weak dual (resp.
of its strong dual) are equicontinuous, and \emph{bornological} if every
set of linear forms on \(E\) that are uniformly bounded on every bounded
subset is equicontinuous. If \(E\) is quasi-complete, then barrelled is
equivalent to quasi-barrelled; in any case, bornological implies
quasi-barrelled.

\section{Topological tensor products}\label{section-1}

\subsection{\texorpdfstring{Generalities on
\(E\hat{\otimes}F\)}{Generalities on E\textbackslash hat\{\textbackslash otimes\}F}}\label{section-1.1}

{[}PTT, chap.~1, §1, no. 1 and no. 3{]}

\medskip
\oldpage{76}

The axiomatic definition of the algebraic tensor product \(E\otimes F\)
of two vector spaces \(E\) and \(F\), and of the canonical bilinear map
\((x,y)\mapsto x\otimes y\) from \(E\times F\) to \(E\otimes F\)
(\autocite{1}) asks only that, for every vector space \(G\), the
bilinear maps from \(E\times F\) to \(G\) correspond bijectively with
linear maps \(f\) from \(E\otimes F\) to \(G\), where the map
corresponding to \(f\) is given by \((x,y)\mapsto f(x\otimes y)\).

\protect\phantomsection\label{section-1-theorem-1}
\begin{itenv}{Theorem 1}
If \(E\) and \(F\) are locally convex spaces, then we can endow
\(E\otimes F\) with a locally convex topology such that, for every
locally convex space \(G\), the \emph{continuous} bilinear maps from
\(E\times F\) to \(G\) correspond exactly to \emph{continuous} linear
maps from \(E\otimes F\) to \(G\). Further, such a topology is unique.

\end{itenv}

Then the \emph{equicontinuous} subsets of \(\mathrm{B}(E,F;G)\)
correspond exactly to the \emph{equicontinuous} subsets of
\(\mathrm{L}(E\otimes F;G)\) as well. Unless otherwise mentioned,
\(E\otimes F\) is assumed to be endowed with the above topology, called
the \emph{projective tensor product of the topologies of \(E\) and
\(F\)}; endowed with this topology, \(E\otimes F\) is called the
\emph{projective topological tensor product} of \(E\) and \(F\).

If \(E\) and \(F\) are normed, then \(E\otimes F\) is normable, and we
can even find a norm such that, for every \emph{normed} space \(G\), the
above isomorphism between \(\mathrm{B}(E,F;G)\) and
\(\mathrm{L}(E\otimes F;G)\) preserves the natural norms. Further, such
a norm is unique. This norm on \(E\otimes F\), denoted by
\(u\mapsto\|u\|\), where the norms of \(E\) and \(F\) are implicit, is
the lower bound of the quantities \(\sum_i\|x_i\|\|y_i\|\) over all
representations of \(u\) in the form \(u=\sum_i x_i\otimes y_i\) (and
this norm has already been considered, in \autocite{8}). It is also the
gauge of the set \(\Gamma(U\otimes V)\), where \(U\) (resp. \(V\)) is
the unit ball in \(E\) (resp. \(F\)), and where \(U\otimes V\) denotes
the set of \(x\otimes y\) such that \(x\in U\) and \(y\in V\) (with
\(\Gamma\) denoting, as per usual, the balanced hull). \oldpage{77} In
the case where \(E\) and \(F\) are general locally convex spaces, a
fundamental system of neighbourhoods of \(0\) in \(E\otimes F\) is
obtained by taking the sets \(\Gamma(U\otimes V)\), where \(U\) (resp.
\(V\)) runs over a fundamental system of neighbourhoods of \(0\) in
\(E\) (resp. \(F\)).

We can introduce the completion of \(E\otimes F\), denoted by
\(E\hat{\otimes}F\), and called the \emph{completed projective tensor
product} of \(E\) and \(F\). If \(E\) and \(F\) are normed spaces, then
\(E\hat{\otimes}F\) is a Banach space (with a well-defined norm!). If
\(E\) and \(F\) are metrisable, then \(E\hat{\otimes}F\) is of type
\((\mathscr{F})\). We have, by definition, the \emph{scholium}: if \(E\)
and \(F\) are locally convex spaces, and \(G\) is a \emph{complete}
locally convex space, then the continuous bilinear maps from
\(E\times F\) to \(G\) correspond bijectively to the continuous linear
maps from \(E\hat{\otimes}F\) to \(G\).

This claim still holds true for equicontinuous sets of maps. In
particular, the dual of \(E\hat{\otimes}F\) is \(\mathrm{B}(E,F)\), with
a correspondence between the equicontinuous subsets (which already
suffices to characterise the induced topology on \(E\otimes F\)).

I do not know if, when \(E\) and \(F\) are of type \((\mathscr{F})\),
this algebraic isomorphism from the dual of \(E\hat{\otimes}F\) to
\(\mathrm{B}(E,F)\) is a \emph{topological} isomorphism, when we endow
\(\mathrm{B}(E,F)\) with the topology given by bibounded convergence,
i.e.~uniform convergence on the products of two bounded sets (``the
problem of topologies''). An equivalent question is the following: is
every bounded subset of \(E\hat{\otimes}F\) contained in the closed
disked hull of a set \(A\hat{\otimes}B\), where \(A\) (resp. \(B\)) is a
bounded subset of \(E\) (resp. \(F\))?

\(\star\) We now give some general tips for calculations involving
\(E\hat{\otimes}F\) ({[}PTT, chap.~1, §1, no. 3{]}). If
\(E=\prod_i E_i\) and \(F=\prod_j F_j\) (as topological-vectorial
products), then \(E\hat{\otimes}F\) can be identified with
\(\prod_{i,j}E_i\hat{\otimes}F_j\). If \(E=\sum_i E_i\) (as a
topological direct sum), and if \(F\) is a normable space, then
\(E\hat{\otimes}F\) can be identified with the topological direct sum
\(\sum_i(E_i\hat{\otimes}F)\). This remains true if \(F\) is an
arbitrary \((\mathscr{DF})\) space, provided that \(I\) is countable,
and these statements can also be generalised to the case where \(E\) is
the inductive limit (in the most general sense) of a family \(E_i\) of
spaces. If \(E\) and \(F\) are both of type \((\mathscr{F})\) (resp.
type \((\mathscr{DF})\)), then so too is \(E\hat{\otimes}F\). Similarly,
if \(E\) and \(F\) are quasi-normable spaces, or Schwartz spaces (see
definitions in \autocite[Section 3]{6}), then so too is
\(E\hat{\otimes}F\). \(\star\)

\subsection{\texorpdfstring{The space \(E\hat{\otimes}F\) when \(E\) and
\(F\) are of type
\((\mathscr{F})\)}{The space E\textbackslash hat\{\textbackslash otimes\}F when E and F are of type (\textbackslash mathscr\{F\})}}\label{section-1.2}

{[}PTT, chap.~1, §2, no. 1{]}

\medskip
\oldpage{78}

\protect\phantomsection\label{section-1-theorem-2}
\begin{itenv}{Theorem 2}
Let \(E\) and \(F\) be \((\mathscr{F})\) spaces. Then every element of
\(E\hat{\otimes}F\) is the sum of an absolutely convergent series of the
form \[
  u = \sum_i \lambda_i x_i \otimes y_i
\] where \((x_i)\) (resp. \((y_i)\)) is a bounded sequence in \(E\)
(resp. \(F\)), and where \((\lambda_i)\) is a summable sequence of
scalars.

\end{itenv}

(It is also true that, if \((x_i)\), \((y_i)\), and \((\lambda_i)\) are
given as above, then the series \(\sum_i\lambda_i x_i\otimes y_i\) is
always absolutely convergent in \(E\hat{\otimes}F\), and so we have a
\emph{characterisation} of the elements of \(E\hat{\otimes}F\)). If
\(E\) and \(F\) are normed, then we can suppose in the above that
\(\|x_i\|\leqslant 1\), \(\|y_i\|\leqslant 1\),
\(\sum_i|\lambda_i|\leqslant\|u\|_1+\varepsilon\), where
\(\varepsilon>0\) is arbitrary and given in advance. In these two
statements, if \(u\) runs over a compact subset of \(E\hat{\otimes}F\),
then we can suppose that the sequences \((x_i)\) and \((y_i)\) remain
fixed (and we can even suppose that they are sequences that tend to
\(0\)), and that \((\lambda_i)\) runs over a compact subset of
\(\ell^1\) (the space of summable sequences). We have an analogous
statement for the concrete representation of convergent sequences in
\(E\hat{\otimes}F\). \hyperref[section-1-theorem-2]{Theorem 2} and its
previous variants serve mainly as a way to:

\begin{itenv}{Corollary}
Let \(E\) and \(F\) be spaces of type \((\mathscr{F})\). Then every
compact subset \(K\) of \(E\hat{\otimes}F\) is contained inside the
canonical image of the unit ball of a space \(E_A\hat{\otimes}F_B\),
where \(A\) (resp. \(B\)) is a compact disked subset of \(E\) (resp.
\(F\)). A fortiori, \(K\) is contained in the closed convex hull of
\(A\otimes B\).

\end{itenv}

This fact also implies that, on \(\mathrm{B}(E,F)\), the ``bicompact
convergence'' topology is identical to the compact convergence topology
on the dual of \(E\hat{\otimes}F\).

\(\star\) For the proof of \hyperref[section-1-theorem-2]{Theorem 2}, we
suppose, for simplicity, that \(E\) and \(F\) are Banach spaces, and let
\(I\) be the product of their unit balls, and \(i\mapsto x_i\) and
\(i\mapsto y_i\) the projections from \(I\) to its factors. It is easy t
see that the linear map
\((\lambda_i)\mapsto\sum_i\lambda_i x_i\otimes y_i\) from \(\ell^1(I)\)
to \(E\hat{\otimes}F\) is a metric homomorphism from the former to a
dense subspace of the latter, and thus \emph{onto} the latter, whence it
indeed follows that \(E\hat{\otimes}F\) can be identified with a
quotient space of \(\ell^1(I)\).

\oldpage{79}

From \hyperref[section-1-theorem-2]{Theorem 2}, we can extract results
of the following type: let \({\mathscr{G}}\) be a locally compact group
(resp. a Lie group), then every summable function \(f\) (resp. every
infinitely differentiable function of compact support) on
\({\mathscr{G}}\) is of the form \(\sum\lambda_i g_i*h_i\), where
\((\lambda_i)\in\ell^1\), and where \((g_i)\) and \((h_i)\) are bounded
sequences in \(\mathrm{L}^i({\mathscr{G}})\) (resp. in
\({\mathscr{D}}({\mathscr{G}})\), the space of infinitely differentiable
functions of compact support on \({\mathscr{G}}\)); we thus immediately
conclude that \(f\) is a linear combination of functions of positive
type that are in \(\mathrm{L}({\mathscr{G}})\) (resp.
\({\mathscr{D}}({\mathscr{G}})\)). In the former case, we can also
restrict to functions that are all of compact support (with the supports
of \(g_i\) and \(h_i\) being contained inside a compact subset that
depends only on the compact support of \(f\)). There is a simple direct
proof in the case of \(\mathrm{L}^1({\mathscr{G}})\), but I do not think
that there is one for \({\mathscr{D}}({\mathscr{G}})\), where the
question presents difficulties even for \({\mathscr{G}}=\mathbb{R}\) (by
using the Fourier transformation). \(\star\)

\subsection{\texorpdfstring{Calculation of
\(\mathrm{L}^1\hat{\otimes}E\)}{Calculation of \textbackslash mathrm\{L\}\^{}1\textbackslash hat\{\textbackslash otimes\}E}}\label{section-1.3}

{[}PTT, chap.~1, §2, no. 2{]}

\medskip

Let \(M\) be a locally compact space endowed with a measure
\(\mu\geqslant 0\), \(E\) a Banach space, and \(\mathrm{L}_E^1(\mu)\)
the space of \(\mu\)-integrable maps from \(M\) to \(E\) (\autocite{2})
endowed with its usual norm: \(\|f\|_1=\int\|f(t)\|\mathrm{d}\mu(t)\).
We denote by \(\mathrm{L}^1(\mu)\) the space of \(\mu\)-summable scalar
functions. Then there exists an obvious bilinear map
\((\varphi,a)\mapsto\varphi\cdot a\) from \(\mathrm{L}^1(\mu)\times E\)
to \(\mathrm{L}_E^1(\mu)\) that is of norm \(\leqslant 1\), and thus
defines a linear map of norm \(\leqslant 1\) from
\(\mathrm{L}^1(\mu)\hat{\otimes}E\) to \(\mathrm{L}_E^1(\mu)\).

\protect\phantomsection\label{section-1-theorem-3}
\begin{itenv}{Theorem 3}
The above map from \(\mathrm{L}^1(\mu)\hat{\otimes}E\) to
\(\mathrm{L}_E^1(\mu)\) is a metric isomorphism from the former space to
the latter.

\end{itenv}

To see this, we can immediately reduce to the case where \(E\) is of
finite dimension, and then proceed by transposition. It then suffices to
apply the classical theorem of Dunford--Pettis that characterises
continuous linear maps from \(\mathrm{L}^1(\mu)\) to \(E'\).

If \(E\) is an arbitrary locally convex space, then we denote by
\(\mathrm{L}_E^1(\mu)\) the completion of the separated space associated
to the space of continuous maps with compact support from \(M\) to
\(E\), endowed with its family of semi-norms
\(f\mapsto\int p(f(t))\mathrm{d}\mu(t)\) (where \(p\) runs over a
fundamental family of continuous semi-norms on \(E\)). Then
\hyperref[section-1-theorem-3]{Theorem 3} easily implies that
\(\mathrm{L}_E^1(\mu)\) is again isomorphic to
\(\mathrm{L}^1(\mu)\hat{\otimes}E\).

\begin{itenv}{Corollary}
If \(E\) is a closed vector subspace of the Banach space \(F\), then the
canonical linear map from \(\mathrm{L}^1(\mu)\hat{\otimes}E\) to
\(\mathrm{L}^1(\mu)\hat{\otimes}F\) is a metric isomorphism.

\end{itenv}

\oldpage{80}

This recovers, for example, the well-known fact that every continuous
linear map from \(E\) to the dual \(\mathrm{L}^\infty(\mu)\) of
\(\mathrm{L}^1(\mu)\) can be extended to a linear map of the same norm
from \(F\) to \(\mathrm{L}^\infty(\mu)\); or, dually, that every
continuous linear map from \(\mathrm{L}^1(\mu)\) to a quotient space
\(F'/E^0\) of a Banach dual by a weakly closed vector subspace comes
from a linear map of the same norm from \(\mathrm{L}^1(\mu)\) to \(F'\).

\(\star\) The analogue of \hyperref[section-1-theorem-3]{Theorem 3} for
\(\mathrm{L}^p\) spaces is false for all \(p>1\). However,
\hyperref[section-1-theorem-3]{Theorem 3} applies, in an essential way,
in many important places in the theory outlined here. We now give some
less important applications (see {[}PTT, chap.~1, §2, no. 2{]} for
details). In \hyperref[section-1-theorem-3]{Theorem 3}, if we take
\(E=c_0\) (the space of scalar sequences that converge to \(0\)), and
note that \(\mathrm{L}_E^1(\mu)\) can then be identified with the space
of latticially bounded sequences in \(\mathrm{L}^1(\mu)\) that converge
to \(0\) almost everywhere, then we see that such sequences in
\(\mathrm{L}^1(\mu)\) form a category of invariant sequences from the
topological-vectorial point of view. In particular, a continuous linear
map from \(\mathrm{L}^1(\mu)\) to a space \(\mathrm{L}^1(\nu)\) sends
latticially bounded sequences that converge to \(0\) almost everywhere
to sequences of the same type. We thus also see that the latticially
bounded subsets of \(\mathrm{L}^1(\mu)\) form an invariant category from
the topological-vectorial point of view. If we take \(F=\ell^p\), with
\(1\leqslant p<+\infty\), then we similarly obtain a
topological-vectorial interpretation of the sequences \((f_i)\) in
\(\mathrm{L}^1(\mu)\) such that \[
  \int \left(
    \sum_i |f_i(t)|^p
  \right)^{1/p}
  \mathrm{d}\mu(t) < +\infty.
\] Such sequences are sent to sequences of the same type by any
continuous linear map from \(\mathrm{L}^1(\mu)\) to a space
\(\mathrm{L}^1(\nu)\). Another interesting application of
\hyperref[section-1-theorem-3]{Theorem 3} is the following: every
bounded subset \(M\) of \(\mathrm{L}^1(\mu)\hat{\otimes}E\) is contained
in the canonical image of the unit ball of a space
\(\mathrm{L}^1(\mu)\hat{\otimes}E_A\), where \(A\) is a closed bounded
disk in the space \(E\) of type \((\mathscr{F})\); a fortiori, \(M\) is
contained in the closed disked hull of \(B\otimes A\), where \(B\) is
the unit ball of \(\mathrm{L}^1(\mu)\), which solves the ``problem of
topologies'' described in \hyperref[section-1.1]{§1.1}. \(\star\)

\subsection{Other examples}\label{section-1.4}

If \(H\) is a Hilbert space, then the elements of \(H'\hat{\otimes}H\),
identified with endomorphisms of \(H\) (the \emph{Fredholm maps} or
\emph{nuclear maps} from \(H\) to \(H\) --- see
\hyperref[section-1.7]{§1.7}) are exactly the endomorphisms \(u\) such
that the positive Hermitian operator \(\sqrt{u^*u}\) is compact and has
a summable sequence of eigenvalues, and \(\|u\|_1\) is then equal to the
sum of the eigenvalues of \(\sqrt{u^*u}\) (repeated with multiplicities,
of course). \oldpage{81} We obtain the operators that have already been
studied in \autocite{4} and \autocite{8}. In fact, \(u\) is also a
Fredholm operator if and only if its Hermitian components
\(\frac12(u+u^*)\) and \(\frac{1}{2i}(u-u^*)\) are, i.e.~if they are
compact Hermitian operators whose sequences of eigenvalues are summable.
\emph{Relation to Hilbert--Schmidt operators: If \(A\) and \(B\) are
Hilbert--Schmidt operators, then \(AB\) is a Fredholm operator, and
\(\|AB\|_1\leqslant\|A\|_2\|B\|_2\)}; and, conversely, every Fredholm
operator \(u\) is the product of two Hilbert--Schmidt operators \(A\)
and \(B\) of norm \(\|A\|_2=\|B\|_2=\sqrt{\|u\|_1}\). All of these facts
are elementary (one we know the spectral decomposition of compact
Hermitian operators to a Hilbert space) and well known.

Numerous other examples of products \(E\hat{\otimes}F\), relating to
\emph{nuclear spaces}, will be seen in \hyperref[section-2.5]{§2.5}.

\(\star\) In the setting of infinite-dimensional Banach spaces, I do now
know, even in particular cases, other concrete characterisations of the
elements of \(E\hat{\otimes}F\) apart from those that we have given.
Also, the elements of \(c_0\hat{\otimes}E\) (where \(E\) can be any
complete locally convex space) can be identified with \emph{certain}
sequences in \(E\) that converge to \(0\), that we can call
\emph{nuclearly convergent to \(0\)}; but if \(E\) is an
infinite-dimensional Banach space, then we always obtain a strictly
smaller class than the class of all sequences that converge to \(0\)
(see \hyperref[section-2-theorem-2]{§2.2, Theorem 2}). We even show that
(if \(E\) is an infinite-dimensional Banach space), for any sequence
\((\lambda_i)\) of positive scalars that is not square-summable, there
exists a sequence \((x_i)\) in \(E\) that does not converge nuclearly to
\(0\), and such that \(\|x_i\|=\lambda_i\) for all \(i\). However, if,
for example, \(E\) is the space \(C(K)\) of continuous functions on a
compact space, then we show that every square-summable sequence in
\(C(K)\) converges nuclearly to \(0\). We also point out that, in any
complete locally convex space, every summable sequence converges
nuclearly to \(0\). \(\star\)

\subsection{\texorpdfstring{\(E\hat{\hat{\otimes}}F\)
spaces}{E\textbackslash hat\{\textbackslash hat\{\textbackslash otimes\}\}F spaces}}\label{section-1.5}

{[}PTT, chap.~1, §3, no. 3{]}

\medskip

If \(E\) and \(F\) are Banach spaces, then \(E\otimes F\) can be
considered as a vector subspace of the Banach space
\(\mathrm{B}(E',F')\) of continuous bilinear forms on \(E'\times F'\).
\oldpage{82} The completion of \(E\otimes F\) under the norm induced by
\(\mathrm{B}(E',F')\) is denoted by \(E\hat{\hat{\otimes}}F\), which is
thus a complete normed vector subspace of \(\mathrm{B}(E',F')\). Every
reasonable normed topology on \(E\otimes F\) is included between the
topology induced by \(E\hat{\otimes}F\) and the topology induced by
\(E\hat{\hat{\otimes}}F\). If \(E\) and \(F\) are now both arbitrary
locally convex spaces, then we can again consider \(E\otimes F\) is a
space of bilinear forms on \(E'\times F'\), and endow it with the
biequicontinuous-convergence topology (i.e.~the topology given by
uniform convergence on the products of an equicontinuous subset of
\(E'\) with an equicontinuous subset of \(F'\)): the completion of
\(E\otimes F\) under this topology is again denoted by
\(E\hat{\hat{\otimes}}F\). If \(E\) and \(F\) are complete, then the
space \(\mathscr{L}_e(E'_s,F'_s)\) of separately continuous bilinear
forms on the product \(E'_s\times F'_s\) of the weak duals \(E'_s\) and
\(F'_s\) endowed with the biequicontinuous-convergence topology is
complete, and so \(E\hat{\hat{\otimes}}F\) can be identified with a
topological vector subspace of \(\mathscr{L}_e(E'_s,F'_s)\). Then the
elements of \(E\hat{\hat{\otimes}}F\) can be identified with certain
separately weakly continuous bilinear maps on \(E'\times F'\), or even
with certain weakly continuous linear maps from \(E'\) to \(F\); these
linear maps send equicontinuous subsets of \(E'\) to relatively compact
subsets of \(F\), and the converse is true in all known cases (see
\hyperref[section-1-appendix-2]{§1, Appendix 2}).

The topology on \(E\otimes F\) induced by \(E\hat{\otimes}F\) is finer
than that induced by \(E\hat{\hat{\otimes}}F\), whence we have a
canonical continuous linear map \[
  E\hat{\otimes}F \to E\hat{\hat{\otimes}}F.
\]

An important, unsolved, problem is the question of whether or not this
map is always bijective (see \hyperref[section-1-appendix-2]{§1,
Appendix 2}). We note that it seems extremely plausible that, if \(E\)
and \(F\) are Banach spaces such that the above map
\(E\hat{\otimes}F\to E\hat{\hat{\otimes}}F\) is a topological
isomorphism (or even simply a topological homomorphism, i.e.~here a map
from the first space \emph{onto} the second), then \(E\) or \(F\) is of
finite dimension. This is true if, for example, \(E\) contains a vector
subspace isomorphic to an \(\ell^p\) space or to \(c_0\).

Let \(E\) and \(F\) be Banach spaces. The norm induced by the dual of
\(E\hat{\otimes}F\) on \(E'\otimes F'\) is clearly the norm induced by
\(E'\hat{\hat{\otimes}}F'\). On the other hand, the norm induced on
\(E'\otimes F'\) by the dual of \(E\hat{\hat{\otimes}}F\) is, in all
known cases (see \hyperref[section-1-appendix-2]{§1, Appendix 2}),
identical to the norm induced by \(E'\hat{\otimes}F'\). This duality,
which barely appears in the present summary, is a precious tool for
various questions (e.g.~{[}PTT, chap.~1, §4, no. 6{]}). \oldpage{83} The
following theorem (which is, truth be told, trivial), can be seen as the
dual response to \hyperref[section-1-theorem-3]{Theorem 3}.

\protect\phantomsection\label{section-1-theorem-4}
\begin{itenv}{Theorem 4}
Let \(M\) be a locally compact space, let \(\mathrm{C}_0(M)\) the space
of continuous scalar functions on \(M\) that are ``zero at infinity'',
endowed with the uniform convergence norm, and let \(E\) be a locally
convex complete space. Then \(\mathrm{C}_0(M)\hat{\hat{\otimes}}E\) is
canonically isomorphic to the space \(\mathrm{C}_0(M,E)\) of continuous
maps from \(M\) to \(E\) that are zero at infinity, endowed with the
uniform convergence topology (and with the uniform norm if \(E\) is a
Banach space). In particular, \(c_0\hat{\hat{\otimes}}E\) is isomorphic
to the space of sequences in \(E\) that tend to \(0\).

\end{itenv}

As with the spaces \(C(K)\hat{\otimes}E\) below, here the spaces
\(\mathrm{L}'(\mu)\hat{\hat{\otimes}}E\) do not, in general, have a
special interpretation as functional spaces. We note, however, that
\emph{the space \(\ell^1\hat{\hat{\otimes}}E\) can be understood as the
space of summable sequences in \(E\)} (i.e.~of ``commutatively
convergent'' sequences in \(E\)).

\subsection{Tensor product of linear maps}\label{section-1.6}

{[}PTT, chap.~1, §1, no. 2{]}

\medskip

Let \(E_i\) and \(F_i\) (for \(i=1,2\)) be locally convex spaces, and
let \(u_i\) be a continuous linear map from \(E_i\) to \(F_i\). In
algebra, we define a linear map \(u_1\otimes u_2\) from
\(E_1\otimes E_2\) to \(F_1\otimes F_2\) by the formula
\((u_1\otimes u_2)(x_1\otimes x_2)=u_1x_1\otimes u_2x_2\). This map is
continuous if we endow \(E_1\otimes E_2\) and \(F_1\otimes F_2\) with
the topologies induced by \(E_1\hat{\otimes}E_2\) and
\(F_1\hat{\otimes}F_2\) (resp. by \(E_1\hat{\hat{\otimes}}E_2\) and
\(F_1\hat{\hat{\otimes}}F_2\)). It then follows that \(u_1\otimes u_2\)
extends, by continuity, to a continuous linear map
\(u_1\hat{\otimes}u_2\) from \(E_1\hat{\otimes}E_2\) to
\(F_1\hat{\otimes}F_2\) (resp. \(u_1\hat{\hat{\otimes}}u_2\) from
\(E_1\hat{\hat{\otimes}}E_2\) to \(F_1\hat{\hat{\otimes}}F_2\)). These
maps will also simply be denoted by \(u_1\otimes u_2\) whenever there is
no fear of confusion.

\protect\phantomsection\label{section-1-theorem-5}
\begin{itenv}{Theorem 5}
If each \(u_i\) is a topological homomorphism from \(E_i\) to a dense
subspace of \(F_i\), then \(u_1\hat{\otimes}u_2\) is a topological
homomorphism from \(E_1\hat{\otimes}E_2\) to a dense subspace of
\(F_1\hat{\otimes}F_2\). If each \(u_i\) is a topological isomorphism
from \(E_i\) to \(F_i\), then \(u_1\hat{\hat{\otimes}}u_2\) is a
topological isomorphism from \(E_1\hat{\hat{\otimes}}E_2\) to
\(F_1\hat{\hat{\otimes}}F_2\).

\end{itenv}

These claims remain true if the \(E_i\) and \(F_i\) are normed and we
consider \emph{metric} homomorphisms and \emph{metric} isomorphisms.
Particularly interesting is the following corollary, which can also be
obtained by an application of \hyperref[section-1-theorem-2]{Theorem 2}.

\begin{itenv}{Corollary}
If the \(E_i\) and \(F_i\) are \((\mathscr{F})\)-spaces (for \(i=1,2\)),
and the \(u_i\) are homomorphisms from \(E_i\) \emph{onto} \(F_i\), then
\(u_1\hat{\otimes}u_2\) is a homomorphism from \(E_1\hat{\otimes}E_2\)
\emph{onto} \(F_1\hat{\otimes}F_2\).

\end{itenv}

\oldpage{84}

As particular cases of this corollary, we obtain interesting lifting
properties of vector functions with values in a quotient space of an
\((\mathscr{F})\)-space \(E\). If e.g.~\(f\) is an infinitely
differentiable map from an open subset of \(\mathbb{R}^n\) to \(E/F\),
then it comes from an infinitely differentiable map from this same open
subset to \(E\). There is an analogous result for holomorphic functions,
or infinitely differentiable functions on \(\mathbb{R}^n\) that decay
rapidly, or summable functions of a certain measure, etc.\footnote{With
  respect to this point, we done that we can prove, by a completely
  different method, the following claim, which can be thought of as dual
  to the corollary of \hyperref[section-1-theorem-3]{Theorem 3}:
  \emph{Let \(M\) be a locally compact and paracompact space (e.g.~a
  compact space), and \(f\) a continuous map from \(M\) to a quotient
  space \(E/F\) of an \((\mathscr{F})\)-space \(E\). Then \(f\) comes
  from a continuous map from \(M\) to \(E\).} Since the space
  \(\mathscr{E}^{(m)}(V)\) of \(m\)-times continuously differentiable
  functions on a infinitely differentiable paracompact manifold \(V\) is
  isomorphic to a direct factor of a space of the form \(C(M)\) (as I
  noted in ``Sur les applications linéaires faiblement compacts
  d'espaces du type \(C(K)\)'', \emph{Can. J. Math.} \textbf{5} (1953),
  p.~144), it thus follows that the analogous lifting theorem holds also
  for \(m\)-times continuously differentiable maps from \(V\) to
  \(E/F\).} Another application is the following: let \(D\) be a
differential operator in the space \(\mathscr{E}(U)\) of infinitely
differentiable functions on an open subset \(U\) of \(\mathbb{R}^n\),
let \(E\) be an \((\mathscr{F})\)-space, and let \(\mathscr{E}(U,E)\) be
the space of infinitely differentiable maps from \(U\) to \(E\). Then we
have
\(\mathscr{E}(U,E)=\mathscr{E}(U)\hat{\otimes}E=\mathscr{E}(U)\hat{\hat{\otimes}}E\)
(see \hyperref[section-2.5]{§2.5}). Let \(D_E\) be the operator in
\(\mathscr{E}(U,E)\) defined by \(D\), and then \(D_E=D\hat{\otimes}1\),
where \(1\) is the identity on \(E\). Then, if \(D\) is a topological
homomorphism (resp. an \emph{onto} topological homomorphism), then so
too is \(D_E\). Indeed, in the case of an onto homomorphism, this is a
particular case of the corollary of
\hyperref[section-1-theorem-5]{Theorem 5}, and, in the general case, we
use \hyperref[section-1-theorem-5]{Theorem 5} as well as the fact that,
for every quotient space \(F\) of \(\mathscr{E}(U)\), we have
\(F\hat{\otimes}E=F\hat{\hat{\otimes}}E\) (since \(\mathscr{E}(U)\), and
thus \(F\), are nuclear; see
\hyperref[section-2-definition-1]{Definition 1} in
\hyperref[section-2.2]{§2.2} and \hyperref[section-2-theorem-3]{Theorem
3} in \hyperref[section-2.3]{§2.3}).

We note that, if \(u_1\) and \(u_2\) are topological isomorphisms, then
\(u_1\hat{\otimes}u_2\) is not, in general, a topological isomorphism
(nor is, in general \(u_1\hat{\hat{\otimes}}u_2\) a topological
homomorphism, if \(u_1\) and \(u_2\) are \emph{onto} topological
homomorphisms). If each \(E_i\) is identified with a
topological-vectorial subspace of \(F_i\) by \(u_i\), then the canonical
map \(u_1\hat{\otimes}u_2\) from \(E_1\hat{\otimes}E_2\) to
\(F_1\hat{\otimes}F_2\) is a topological isomorphism if and only if
every equicontinuous set of bilinear forms on \(E_1\times E_2\) is the
set of restrictions of an equicontinuous set of bilinear forms on
\(F_1\times F_2\). \oldpage{85} If \(F_1\) and \(F_2\) are of type
\((\mathscr{F})\), then it suffices to consider sets consisting of
\emph{one single} bilinear form. In general, this criterion will not
hold true, but it is linked to an existence problem of topological
complements.

\(\star\) More precisely, if \(E_1\) and \(E_2\) are direct factors,
then \(E_1\hat{\otimes}E_2\) can be identified with a
topological-vectorial subspace of \(F_1\hat{\otimes}F_2\). Also, if
\(E\) is a topological-vectorial subspace of the Banach space \(F\), and
if the canonical map \(E\hat{\otimes}G\to F\hat{\otimes}G\) is a
topological isomorphism when \(G=F'\), then \(E''\) is a direct factor
of \(F''\); thus, in the frequent case where \(E\) is already a direct
factor of \(E''\), \(E\) will also be a direct factor of \(F\).

A useful case where the tensor product \(u\hat{\otimes}v\) of two
topological isomorphisms is a topological isomorphism is the following:
\emph{If \(E''\) is the bidual of \(E\), then \(E\hat{\otimes}F\) can be
identified with a topological-vectorial subspace (resp. a normed vector
subspace, if \(E\) and \(F\) are normed) of \(E''\hat{\otimes}F\)}.
\(\star\)

\subsection{Nuclear maps}\label{section-1.7}

{[}PTT, chap.~1, §3, no. 2{]} \medskip

Let \(E\) and \(F\) be Banach spaces. The continuous bilinear map
\((x',y)\mapsto x'\otimes y\) from \(E'\times F\) to \(\mathrm{L}(E,F)\)
defines a natural continuous linear map from \(E'\hat{\otimes}F\) to
\(\mathrm{L}(E,F)\); the elements of the image of \(E'\hat{\otimes}F\)
in \(\mathrm{L}(E,F)\) are said to be \emph{nuclear maps} from \(E\) to
\(F\). (We also define, between arbitrary locally convex spaces, the
notions of \emph{trace maps} and \emph{Fredholm maps} --- see
\hyperref[section-1-appendix-1]{Appendix 1} --- which coincide with the
notion of nuclear maps if \(E\) and \(F\) are Banach spaces; in this
case, we can thus freely switch between saying ``nuclear maps'', ``trace
maps'', or ``Fredholm maps''.) Nuclear maps from \(E\) to \(F\) form a
vector space which can be identified with a quotient space of
\(E'\hat{\otimes}F\) (and with \(E'\hat{\otimes}F\) itself in all known
cases --- see ``problem of bijectivity'' in \hyperref[section-1]{§1}).
The quotient norm, again denoted by \(u\mapsto\|u\|_1\), is called the
\emph{trace norm} of the nuclear operator \(u\).

If \(E\) and \(F\) are arbitrary locally convex spaces, then a linear
map \(u\) from \(E\) to \(F\) is said to be \emph{nuclear} if it is the
composition of a sequence of three operators \[
  E \xrightarrow{\alpha} E_1 \xrightarrow{\beta} F_1 \xrightarrow{\gamma} F
\] where \(E_1\) and \(F_1\) are Banach spaces, \(\beta\) is a nuclear
map from \(E_1\) to \(F_1\), and \(\alpha\) and \(\gamma\) are
continuous linear maps. \oldpage{86} It is equivalent to say that there
exists an weakly closed, equicontinuous, disked subset \(A\) of \(E'\),
and a bounded disked subset \(B\) of \(F\) such that \(F_B\) is
complete, and finally some \(u_0\in E'_A\hat{\otimes}F_B\) such that
\(u\) is the operator from \(E\) to \(F\) defined by \(u_0\) (\emph{a
priori}, \(u_0\) defines a nuclear map from \(\widehat{E}_{A^0}\) to
\(F_B\)). \emph{A nuclear map is always compact} (i.e.~sends a suitable
neighbourhood of \(0\) to a relatively compact set). Even better: by the
corollary of \hyperref[section-1-theorem-2]{Theorem 2}, we can suppose
(in the above) that \(B\) and \(A\) are \emph{compact} subsets of \(F\)
and strong \(E'\) (respectively). \hyperref[section-1-theorem-2]{Theorem
2}, applied directly, also gives the following: \emph{nuclear maps from
\(E\) to \(F\) are exactly the maps that are sums of series (always
absolutely convergent in \(\mathrm{L}(E,F)\) endowed with the
bounded-convergence topology) \(u=\sum\lambda_i x'_i\otimes y_i\), where
\((x'_i)\) is an equicontinuous sequence in \(E'\), \((y_i)\) is a
sequence contained in a compact disk of \(F\), and \((\lambda_i)\) is a
summable sequence of scalars}. If we compose a nuclear map, on the left
or on the right, with a continuous linear map, then we obtain another
nuclear map. The dual of a nuclear map from \(E\) to \(F\) is a nuclear
map from strong \(F'\) to strong \(E'\) (and even from \(F'\) endowed
with the uniform convergence on compact disks of \(F\), to strong
\(E'\)).

Nuclear maps from \(E\) to itself (more precisely, the slightly larger
category of Fredholm maps from \(E\) to itself) form a natural domain
for \emph{Fredholm theory}. Here, our interest lies in other properties
of these operators, that result directly from either
\hyperref[section-1-theorem-2]{Theorem 2} or the corollary of
\hyperref[section-1-theorem-5]{Theorem 5}.

\protect\phantomsection\label{section-1-theorem-6}
\begin{itenv}{Theorem 6}

Let \(E\) and \(G\) be locally convex spaces, and \(F\) a vector
subspace of \(E\). Then:

\begin{enumerate}
\def\labelenumi{\alph{enumi}.}
\item
  Every nuclear map from \(F\) to \(G\) is the restriction of a nuclear
  map from \(E\) to \(G\).
\item
  Suppose that \(F\) is closed, and that every compact disc of \(E/F\)
  is contained in the canonical image of a bounded disc \(A\) of \(E\)
  such that \(E_A\) is complete (for example, a \emph{complete} bounded
  disc). (It suffices, for example, for \(E\) to be an
  \((\mathscr{F})\)-space, or for it be the dual of an
  \((\mathscr{F})\)-space and for \(F\) to be \emph{weakly} closed.)
  Then every nuclear map from \(G\) to \(E/F\) can be obtained from a
  nuclear map from \(G\) to \(E\) by passing to the quotient.
\end{enumerate}

\end{itenv}

We have analogous statements for ``equinuclear'' sets of maps, if by
that we mean a set of maps from a locally convex space \(M\) to another
locally convex space \(N\), contained in the set of maps defined by the
unit ball of a space \(M'_A\hat{\otimes}N_B\), where \(A\) is a weakly
closed equicontinuous disk of \(M'\), and \(B\) a bounded disk in \(N\)
such that \(N_B\) is complete.

\oldpage{87}

We note that, if \(u\) is a nuclear map from \(E\) to \(F\), \(M\) a
closed vector subspace of \(E\) contained in the kernel of \(u\), and
\(N\) a closed vector subspace of \(F\) containing \(u(E)\), then, in
general, the linear map from \(E/M\) to \(F\) or from \(E\) to \(N\)
defined by \(u\) \emph{is not nuclear}, even if \(E\) and \(F\) are
reflexive Banach spaces. However, if \(M\) (resp. \(N\)) admits a
topological complement, then the map from \(E/M\) to \(F\) (resp. from
\(E\) to \(N\)) defined by \(u\) will itself be nuclear. This is the
case, in particular, if \(E\) (resp. \(F\)) is a Hilbert space.

\subsection{Integral linear maps, integral bilinear
forms}\label{section-1.8}

{[}PTT, chap.~1, §4, no. 3 and no. 4{]} \medskip

\protect\phantomsection\label{section-1-theorem-7}
\begin{itenv}{Theorem 7}

Let \(E\) and \(F\) be locally convex (resp. normed) spaces, and \(v\) a
separately continuous bilinear form on \(E\times F\). Then the following
conditions are equivalent:

\begin{enumerate}
\def\labelenumi{\alph{enumi}.}
\item
  \(u\mapsto\langle u,v\rangle\) is a linear form on \(E\otimes F\) that
  is continuous for the topology induced by \(E\hat{\hat{\otimes}}F\)
  (resp. \(u\mapsto\langle u,v\rangle\) is a linear form on
  \(E\otimes F\) of norm \(\leqslant 1\) when \(E\otimes F\) is endowed
  with the norm induced by \(E\hat{\hat{\otimes}}F\)).
\item
  \(v\) is contained in the closed disked hull in
  \(\mathscr{B}_\mathrm{s}(E,F)\) (the space \(\mathscr{B}(E,F)\)
  endowed with the simple-convergence topology) of a set \(M\otimes N\),
  where \(M\) is an equicontinuous subset of \(E'\), and \(N\) an
  equicontinuous subset of \(F'\) (resp. \(M\) the unit ball of \(E'\),
  and \(N\) the unit ball of \(F'\)).
\item
  There exists a measure \(\mu\) on the product space of a weakly
  compact equicontinuous subspace \(M\) of \(E'\) with a weakly compact
  equicontinuous subspace \(N\) of \(F'\) (resp. a measure \(\mu\) of
  norm \(\leqslant 1\) on the product of the unit ball of \(E'\) with
  the unit ball of \(F'\), endowed with the product of the weak
  topologies) such that we have the formula \[
     v = \int_{M\times N} x'\otimes y' \mathrm{d}\mu(x',y')
   \] (the weak integral in \(\mathscr{B}(E,F)\), in duality with
  \(E\otimes F\)).
\item
  There exists a compact space endowed with a positive measure \(\mu\)
  of norm \(\leqslant 1\), a continuous linear map \(\alpha\) from \(E\)
  to \(\mathrm{L}^\infty(\mu)\), and a continuous linear map \(\beta\)
  from \(F\) to \(\mathrm{L}^\infty(\mu)\) (resp. the same, but with
  \(\alpha\) and \(\beta\) also of norm \(\leqslant 1\)) such that we
  have \(u(x,y)=\langle\alpha x,\beta y\rangle\) for \(x\in E\),
  \(y\in F\).
\end{enumerate}

\end{itenv}

\oldpage{88}

A bilinear form on \(E\times F\) is said to be \emph{integral} if it
satisfies any of the equivalent conditions of
\hyperref[section-1-theorem-7]{Theorem 7}. In particular, the dual of
\(E\hat{\hat{\otimes}}F\) can be identified with the space
\(\mathrm{J}(E,F)\) of integral bilinear forms on \(E\times F\). If
\(E\) and \(F\) are Banach spaces, then \(\mathrm{J}(E,F)\) will be
endowed with the dual norm of the Banach space
\(E\hat{\hat{\otimes}}F\), which we call the \emph{integral norm}, and
denote by \(\|v\|'_1\). Similarly, a linear map \(v\) from one locally
convex space \(E\) to another \(G\) is said to be \emph{integral} if the
corresponding bilinear form on \(E\times G'\) is integral.

If \(E\) and \(G\) are Banach spaces, then we also call the integral
norm of the corresponding bilinear form of \(v\) the integral norm of
\(v\).

Recall that, in all known cases, when \(E\) and \(F\) are Banach, the
natural linear map from \(E'\hat{\otimes}F'\) to \(\mathrm{J}(E,F)\) is
a metric isomorphism from the first space to the second (see
\hyperref[section-1.5]{§1.5}), which is why we use the notation
\(\|v\|'_1\) for the integral norm, closely related to the notation
\(\|v\|_1\) for the trace norm. Criterion (d) of
\hyperref[section-1-theorem-7]{Theorem 7} takes the following form
(which we state for Banach spaces, for simplicity) for integral linear
maps: \emph{Let \(v\) be a linear map from one Banach space \(E\) to
another \(F\). Then \(v\) is integral and of integral norm
\(\leqslant 1\) if and only if the map from \(E\) to \(F''\) that it
defines can be obtained by composing a linear map of norm \(\leqslant\)
from \(E\) to some space \(\mathrm{L}^\infty(\mu)\), constructed with a
suitable positive measure of norm \(\leqslant 1\) on a compact space,
with the identity map from \(\mathrm{L}^\infty(\mu)\) to
\(\mathrm{L}^1(\mu)\), and finally with a linear map of norm
\(\leqslant\) from \(\mathrm{L}^1(\mu)\) to \(F''\).} Similarly,
criterion (b) of \hyperref[section-1-theorem-7]{Theorem 7} easily gives:
\emph{The linear map \(v\) from \(E\) to \(F\) is integral and of
integral norm \(\leqslant 1\) if and only if it is an adherent point in
\(\mathrm{L}_s(E,F_s)\) (where \(F_s\) denotes \(F\) endowed with the
weak topology, and \(\mathrm{L}_s(E,F_s)\) denotes \(\mathrm{L}(E,F_s)\)
endowed with the simple-convergence topology) of the disked hull of the
set of the \(x'\otimes y\), where \(x'\) (resp. \(y\)) runs over the
unit ball of \(E'\) (resp. of \(F\)); or if it is adherent to the set of
nuclear operators of trace-norm \(\leqslant 1\).}

\begin{rmenv}{Examples}
Let \(E\) and \(F\) be arbitrary locally convex spaces. Then every
bilinear form on \(E\times F\) defined by an element of a space
\(E'_A\hat{\otimes}F'_B\), where \(A\) (resp. \(B\)) is a weakly closed
disked equicontinuous subset of \(E'\) (resp. of \(F'\)), is integral.
Thus every nuclear map from \(E\) to \(F\) is integral. \oldpage{89} The
converse is false, even for Banach spaces, since the identity map from
\(\mathrm{L}^\infty(\mu)\) to \(\mathrm{L}^1(\mu)\) is integral, but it
is not, in general, even compact. If \(E=\mathrm{C}(M)\) (resp.
\(F=\mathrm{C}(N)\)) is the space of continuous scalar functions on the
compact space \(M\) (resp. \(N\)), then we have seen
(\hyperref[section-1-theorem-4]{Theorem 4}) that
\(\mathrm{C}(M)\hat{\otimes}\mathrm{C}(N)\) can be identified with its
norm to the space \(\mathrm{C}(M\times N)\), and so the space of
integral bilinear forms on \(\mathrm{C}(M)\times\mathrm{C}(N)\) can be
identified with its norm to the space of Radon measures on the compact
space \(M\times N\). Other examples will be seen in
\hyperref[section-1.9]{§1.9}.

By composing an integral linear map on the left or on the right with a
continuous linear map, we obtain another integral linear map. The
transpose of an integral map from \(E\) to \(F\) is an integral map from
strong \(F'\) to strong \(E'\).

Using, for example, criterion (a) of
\hyperref[section-1-theorem-7]{Theorem 7}, we see that, if \(E_1\) and
\(E_2\) are locally convex spaces, and \(F_1\) (resp. \(F_2\)) is a
topological vector subspace of \(E_1\) (resp. of \(E_2\)), then every
integral bilinear form on \(F_1\times F_2\) can be extended to a
integral bilinear form on \(E_1\times E_2\), with equal integral norm if
\(E_1\) and \(E_2\) are normed (compare with
\hyperref[section-1-theorem-6]{Theorem 6}). The most important
properties of integral maps are summarised in the following:

\end{rmenv}

\protect\phantomsection\label{section-1-theorem-8}
\begin{itenv}{Theorem 8}

Let \(u\) be an integral linear map from a locally convex space \(E\) to
a locally convex space \(F\).

\begin{enumerate}
\def\labelenumi{\arabic{enumi}.}
\item
  If \(F\) is quasi-complete, then \(u\) is weakly compact, and sends
  weakly compact subsets of \(E\) to compact subsets of \(F\). If \(v\)
  is a linear map from \(F\) to a locally convex space \(G\) that sends
  bounded subsets to weakly relatively compact subsets, then
  \(v\circ u\) is a compact map.
\item
  Let \(v\) be a linear map from \(F\) to an \((\mathscr{F})\)-space
  \(G\) that sends bounded subsets to weakly relatively compact subsets
  (resp. a linear map from a \((\mathscr{DF})\)-space \(G\) to \(E\)
  that sends bounded subsets to weakly relatively compact subsets of
  \(E\)). Then \(v\circ u\) (resp. \(u\circ v\)) is a \emph{nuclear} map
  from \(E\) to \(G\) (resp. from \(G\) to \(F''\)). If \(E\), \(F\),
  and \(G\) are Banach spaces, then
  \(\|v\circ u\|'_1\leqslant\|v\|\|u\|'\).
\end{enumerate}

\end{itenv}

\protect\phantomsection\label{section-1-corollary-1}
\begin{itenv}{Corollary 1}
The composition of two integral maps is a nuclear map.

\end{itenv}

\begin{rmenv}{Other corollaries}
An integral map from \(E\) to \(F\) is nuclear if \(F\) is a reflexive
\((\mathscr{F})\)-space, and it is a nuclear map from \(E\) to \(F''\)
if \(E\) is a reflexive \((\mathscr{DF})\)-space. \oldpage{90} An
integral map from \(E\) to \(F\) sends summable sequences to absolutely
summable sequences, and weakly convergent sequences to nuclearly
convergent sequences (see the end of \hyperref[section-1.4]{§1.4}).

\end{rmenv}

\(\star\) So, considering the unit circumference \(T\) of the complex
plane, endowed with its Haar measure \(\mu\), suppose that the sequence
\((a_n)\) in the set \(\mathbb{Z}\) of integers is such that
\((\varepsilon_n a_n)\) is the sequence of coefficients of a function in
\(\mathrm{L}^\infty(\mu)\), for some sequence \((\varepsilon)\) of
numbers equal to \(+1\) or \(-1\); then we can easily see that the
sequence \(a_nz^n\in\mathrm{L}^\infty(\mu)\) is summable, and is thus an
absolutely summable sequence in \(\mathrm{L}^1(\mu)\), whence,
immediately, \((a_n)\in\ell^1\). We have thus obtained an analogue of
the well-known theorem of Littlewood (in which \(\mathrm{L}^\infty\) is
replaced by \(\mathrm{L}^1\), and \(\ell^1\) by \(\ell^2\)). We note
that Littlewood's theorem can be obtained in the same way, as a
consequence of the following theorem, which has much more general
consequences (and which will later be published, along with some of its
various consequences): \emph{every summable sequence in an
\(\mathrm{L}^1(\mu)\)-space (for an arbitrary measure) has a sequence of
norms that is square summable} (and even belonging to
\(\ell^2\hat{\otimes}\mathrm{L}^1(\mu)\); this is dual to the theorem,
noted at the end of \hyperref[section-1.4]{§1.4}, that says that every
square summable sequence in a \(\mathrm{C}(K)\)-space --- and even every
sequence belonging to \(\ell^2\hat{\hat{\otimes}}\mathrm{C}(K)\) --- is
nuclearly convergent to \(0\)).

We will give some hints concerning the proof of
\hyperref[section-1-theorem-8]{Theorem 8}, that rely in an essential
manner on criterion (d) of \hyperref[section-1-theorem-7]{Theorem 7}.
The fact that \(u\) is a weakly compact map follows immediately from the
fact that the identity map from \(\mathrm{L}^\infty(\mu)\) to
\(\mathrm{L}^1(\mu)\) is a weakly compact map. The other claims of the
theorem follow easily from the second part, which is more difficult. We
can reduce to showing that every weakly compact linear map from
\(\mathrm{L}^1(\mu)\) to an \((\mathscr{F})\)-space \(G\) induces a
nuclear map from \(\mathrm{L}^\infty(\mu)\) to \(G\), and thus
(\hyperref[section-1-theorem-3]{Theorem 3}) comes from some
\(f\in\mathrm{L}^1_G(\mu)\). But a theorem of Dunford--Pettis--Phillips
tells us that such a map is in fact given by a \emph{strongly measurable
and bounded} map \(f\) from \(M\) to \(G\). Note that
\hyperref[section-1-corollary-1]{Corollary 1}, which is important due to
its application in the theory of nuclear spaces, admits a simpler direct
proof: we can reduce to showing that, if \(\mu\) and \(\nu\) are
measures on compact \(M\) and \(N\), then a continuous linear map \(u\)
from \(\mathrm{L}^1(\mu)\) to \(\mathrm{L}^\infty(\nu)\) defines a
\emph{nuclear} map from \(\mathrm{L}^\infty(\mu)\) to
\(\mathrm{L}^1(\nu)\). But \(u\) can be identified with a continuous
linear form on \(\mathrm{L}^1(\mu)\hat{\otimes}\mathrm{L}^1(\nu)\),
which is isomorphic to the space \(\mathrm{L}^1(\mu\otimes\nu)\)
(\hyperref[section-1-theorem-3]{Theorem 3}), and so \(u\) is defined by
a measurable and bounded function \(f\) on \(M\times N\), which can be
identified a fortiori with an element of
\(\mathrm{L}^1(\mu\otimes\nu)=\mathrm{L}^1(\mu)\hat{\otimes}\mathrm{L}^1(\nu)\).
This latter element indeed defines a nuclear map from
\(\mathrm{L}^\infty(\mu)\) to \(\mathrm{L}^1(\nu)\), which happens to be
exactly the map induced by \(u\). QED.

\oldpage{91}

In PTT, we deduce \hyperref[section-1-theorem-8]{Theorem 8} from more
general results (see {[}PTT, chap.~1, §4, no. 2{]}). The majority of
{[}PTT, chap.~1, §4{]} (the densest part of the whole work) is dedicated
to the exposition of these results and their many consequences, which we
cannot give in this summary. \(\star\)

\subsection{\texorpdfstring{Integral linear maps to an \(\mathrm{L}^1\)
space or a \(\mathrm{C}_0(M)\)
space}{Integral linear maps to an \textbackslash mathrm\{L\}\^{}1 space or a \textbackslash mathrm\{C\}\_0(M) space}}\label{section-1.9}

{[}PTT, chap.~1, §4, no. 4{]}

\medskip

\(\star\) We can characterise the integral linear from a locally convex
space \(E\) to an \(\mathrm{L}^1(\mu)\) space (where \(\mu\) is an
arbitrary measure on a locally compact space \(M\)): they are the linear
maps that send a suitable neighbourhood \(V\) of \(0\) in \(E\) to a
\emph{latticially bounded} subset of \(\mathrm{L}^1(\mu)\). If \(E\) is
normed, with \(V\) its unit ball, and \(h=\sup_{x\in V}|ux|\) (so that
\(h\) is a positive element of \(\mathrm{L}^1(\mu)\)), then
\(\|u\|'_1=\|h\|_1\). If \(E\) or \(\mathrm{L}^1(\mu)\) is separable,
then the theorem of Dunford--Pettis gives an equivalent criterion (which
we state, as an example, in the case where \(E\) is assumed to be
normed): \emph{there exists a weakly measurable map \(f\) from \(M\) to
\(E'\) such that \(\|f(t)\|\) is a summable function of \(t\), and such
that, for all \(x\in E\), \(ux\) is the class in \(\mathrm{L}^1(\mu)\)
of the function \(t\mapsto\langle x,f(t)\rangle\); then
\(\|u\|'_1\leqslant\int\|f(t)\|\mathrm{d}\mu(t)\) (and we have equality
for a suitable choice of \(f\)).} We can also characterise the
\emph{nuclear} maps from a locally convex space \(E\) to
\(\mathrm{L}^1(\mu)\): they are those which send a suitable
neighbourhood \(V\) of \(0\) in \(E\) to a subset \(A\) of
\(\mathrm{L}^1(\mu)\) which is latticially bounded and further
\emph{equimeasurable} (by which we mean that, for every compact
\(K\subset M\), and every \(\varepsilon>0\), there exists a compact
\(K_0\) such that \(\mu(K\cap K_0^c)\leqslant\varepsilon\), and such
that the \(\varphi\in A\) agree almost everywhere on \(K_0\) with the
functions from an equicontinuous and uniformly bounded set of functions
on \(K_0\)). If we suppose, for simplicity, that \(E\) is a Banach
space, then this is equivalent to saying that the map in question is
given by an \emph{integrable} \autocite{2} map \(f\) from \(M\) to
\(E'\) (and, indeed, it follows immediately from
\hyperref[section-1-theorem-3]{Theorem 3} that this indeed implies that
\(u\) is nuclear).

\oldpage{92}

Dually, let \(M\) be a locally compact space, and \(E\) a Banach space
(for simplicity), and suppose that either \(M\) is metrisable and
countable at infinity, or that \(E\) is separable. Then the integral
maps \(u\) from \(\mathrm{C}_0(M)\) (the space of continuous functions
on \(M\) that are ``zero at infinity'') to \(E'\), i.e.~the continuous
linear forms on
\(\mathrm{C}_0(M,E)=\mathrm{C}_0(M)\hat{\hat{\otimes}}E\), are exactly
those given by a measure \(\mu\) on \(M\) and a weakly
\(\mu\)-measurable and bounded map \(f\) from \(M\) to \(E'\) by the
formula \(u\varphi=\int\varphi(t)f(t)\mathrm{d}\mu(t)\). We can further
suppose that \(\|f(t)\|=1\) for all \(t\), and that
\(\|\mu\|_1=\|u\|'_1\). Using \hyperref[section-1-theorem-3]{Theorem 3},
we also find that the \emph{nuclear} maps from \(\mathrm{C}_0(M)\) to
\(E'\), or even from \(\mathrm{C}_0(M)\) to an arbitrary Banach space
\(F\), are exactly those given by a pair \((u,f)\) as above, but with
\(f\) being an \emph{integrable} map from \(M\) to \(F\); it is no
longer necessary here to make any separability hypotheses. In
particular, if \(K\) and \(L\) are compact spaces, \(\mu\) a measure on
\(K\), and \(N(x,y)\) a continuous scalar function on \(K\times L\),
then the map \(f\mapsto\int f(x)N(x,y)\mathrm{d}\mu(x)\) from \(C(K)\)
to \(C(L)\), defined by the continuous kernel \(N\), is nuclear. Note
that we have just described two remarkable categories of ``vectorial
measures'' on \(M\), which we call \emph{integral vectorial measures}
and \emph{nuclear vectorial measures} (respectively) on \(M\).

The fact that we can characterise the integral and nuclear maps from a
Banach space (for example) \(E\) to \(\mathrm{L}^1(\mu)\) by properties
concerning their images of the unit ball is completely special to
\(\mathrm{L}^1(\mu)\) (see the end of \hyperref[section-1.7]{§1.7}) and
is also linked to the corollary of
\hyperref[section-1-theorem-3]{Theorem 3}. Similarly, we can show that
the integral linear maps from \(\mathrm{L}^1(\mu)\) to a locally convex
space \(E\) can be characterised by properties concerning their images
of the unit ball of \(\mathrm{L}^1(\mu)\), a feature that we will not
further study here (see {[}PTT, chap.~1, §4, no. 6{]}). \(\star\)

\subsection*{Appendix 1: Various variants of the notion of topological
tensor product}\label{section-1-appendix-1}
\addcontentsline{toc}{subsection}{Appendix 1: Various variants of the
notion of topological tensor product}

{[}PTT, chap.~1, §3{]}

On \(E\otimes F\), we can introduce a large number of distinct
interesting topologies (even when \(E\) and \(F\) are Banach spaces). We
will content ourselves here with noting the existence of a unique
locally convex topology on \(E\otimes F\) such that, for every locally
convex space \(G\), \emph{separately continuous} bilinear maps from
\(E\times F\) to \(G\) correspond exactly to \emph{continuous} linear
maps from \(E\otimes F\) to \(G\). To separately equicontinuous sets of
bilinear maps from \(E\times F\) then correspond equicontinuous sets of
linear maps from \(E\otimes F\). \oldpage{93} Endowed with this
topology, \(E\otimes F\) is called the \emph{inductive tensor product}
of \(E\) and \(F\), and its completion, denoted by \(E\bar{\otimes}F\).
Its dual is thus the space \(\mathscr{B}(E,F)\), and the equicontinuous
subsets of this dual are the separately equicontinuous sets of bilinear
forms on \(E\times F\) (which already suffice to determine the topology
of the inductive tensor product). The inductive tensor product topology
on \(E\otimes F\) is finer than the projective tensor product topology,
and these two topologies are identical if and only if the separately
equicontinuous sets of bilinear forms on \(E\times F\) are already
equicontinuous (for example, if \(E\) and \(F\) are of type
\((\mathscr{F})\), or if \(E\) and \(F\) are of type \((\mathscr{DF})\)
and barrelled). So if \(E\) is a non-normable space, then the two
topologies above on \(E\otimes E'\) give distinct duals (since the
canonical bilinear form on \(E\times E'\) is separately continuous but
not continuous). --- If \(E\) is the inductive limit (in the general
sense) of a family of spaces \(E_i\), and \(F\) the inductive limit of a
family of spaces \(F_j\), then the topological-vectorial subspace \(H\)
of \(E\bar{\otimes}F\) generated by the canonical images of the spaces
\(E_i\bar{\otimes}F_j\) is the inductive limit of these latter spaces
(whence the name ``inductive'' tensor product). Unfortunately, the space
\(H\) above is often not complete, i.e.~it is distinct from
\(E\bar{\otimes}F\). --- The analogous statement of the above, for when
\(E\) and \(F\) are \emph{product} spaces, is only true under fairly
restrictive conditions, e.g.~if \(F\) is the topological-vectorial
product of a family of spaces \emph{of type \((\mathscr{F})\)}. ---
Finally, note that the notion of topological tensor product that we have
just developed gives rise to a notion of tensor product of two
continuous linear maps, completely analogous to the notion developed in
\hyperref[section-1.6]{§1.6}.

We should define a remarkable dense subset of \(E\bar{\otimes}F\), which
is more important that \(E\bar{\otimes}F\) itself (even though it is
often not complete, since it is distinct from \(E\bar{\otimes}F\)): it
is the subspace given by the union of the canonical images of the spaces
\(E_A\bar{\otimes}F_B=E_A\hat{\otimes}F_B\), where \(A\) (resp. \(B\))
is a bounded disk of \(E\) (resp. \(F\)) such that \(E_A\) (resp.
\(F_B\)) is complete. The elements of this subspace are called
\emph{Fredholm kernels} in \(E\bar{\otimes}F\), and they appear, for
example, in the theory of nuclear spaces (see
\hyperref[section-2-corollary-3]{Corollary 3 of Theorem 1 in §2.2}). The
subspaces of Fredholm kernels are sent to one another by tensor products
of continuous linear maps. --- \hyperref[section-1-theorem-2]{Theorem 2
(in §1.2)} shows that, if \(E\) and \(F\) are of type \((\mathscr{F})\),
then every element of \(E\bar{\otimes}F=E\hat{\otimes}F\) is already a
Fredholm kernel; it further gives an explicit structure theorem for
Fredholm kernels in the general case: \oldpage{94} Fredholm kernels in
\(E\bar{\otimes}F\) can be represented by series \[
  u = \sum\lambda_i x_i\otimes y_i
\] where \((x_i)\) (resp. \((y_i)\)) is a series extracted from a
compact disk of \(E\) (resp. \(F\)), and where \((\lambda_i)\) is a
summable sequence of scalars. Thus \(u\) itself comes from an element of
some space \(E_A\hat{\otimes}F_B\), where \(A\) and \(B\) are
\emph{compact} disks.

We define a \emph{Fredholm map} from \(E\) to \(F\) to be a map defined
by a Fredholm kernel of \(E'\bar{\otimes}F\). Such a map is weakly
continuous, but not necessarily continuous; but if every strongly
compact disk of \(E'\) is equicontinuous (in particular, if the topology
on \(E\) is the Mackey topology \(\tau(E,E')\)), then every Fredholm map
from \(E\) to \(F\) is already nuclear, and a fortiori continuous. ---
Note that it is the Fredholm kernels of \(E'\bar{\otimes}E\) that form
the natural domain for Fredholm theory.

Let \(E\) be a locally convex space, and take a locally convex topology
on its dual that is less fine than the weak topology. Then the canonical
bilinear form on \(E\times E'\) is separately continuous, and thus
defines a continuous linear form on \(E\bar{\otimes}E'\), called the
\emph{trace} form. Let \(F\) be another locally convex space; let
\(A\in\mathscr{B}(E,F)\), and suppose that the corresponding linear map
\({}^tA\) from \(F\) to \(E'\) is continuous (which will be the case,
for example, if we take the strong topology on \(E'\) and assume that
\(F\) is barrelled); then \(1\otimes{}^tA\) is a continuous linear map
from \(E\bar{\otimes}F\) to \(E\bar{\otimes}E'\), also denoted by
\(u\mapsto {}^tAu\). Then the passage to the trivial limit gives, for
all \(u\in E\bar{\otimes}F\), \[
  \langle u,A\rangle = \operatorname{Tr}{}^tAu.
\]

This makes the duality between \(E\bar{\otimes}F\) and
\(\mathscr{B}(E,F)\) explicit by means of the trace form. This formula
immediately generalises for the natural pairing corresponding to any
``reasonable'' type of completed topological tensor product (PTT, §3,
no. 3, prop. 17). --- If \(K\) is a compact space endowed with a measure
\(\mu\), then we have already seen (cf.~\hyperref[section-1.9]{§1.9})
that the integral operator defined by a continuous kernel \(N(x,y)\)
(defined on \(K\times K\)) is a nuclear operator in \(\mathrm{C}(K)\).
We can easily show that its trace is exactly
\(\int N(x,x)\mathrm{d}\mu(x)\).

\subsection*{Appendix 2: The properties and problems of
approximation}\label{section-1-appendix-2}
\addcontentsline{toc}{subsection}{Appendix 2: The properties and
problems of approximation}

{[}PTT, chap.~1, §5{]}

The most important problem that remains to be solved in the theory of
topological tensor products is the following ``\emph{bijectivity
problem}'': \oldpage{95} is the canonical map from \(E\hat{\otimes}F\)
to \(E\hat{\hat{\otimes}}F\) always bijective? By the Hahn--Banach
theorem, this problem can be turned into of the variations of the
``\emph{approximation problem}'': is every continuous bilinear form on
\(E\times F\) the limit, in the weak topology defined by
\(E\hat{\otimes}F\), of degenerate continuous bilinear forms, i.e.~of
those coming from \(E'\otimes F'\)? Under this form of the problem, we
see that we can reduce to the case where \(E\) and \(F\) are Banach
spaces. Since then the bicompact-convergence topology (i.e.~uniform
convergence on products of a compact of \(E\) with a compact of \(F\))
on \(\mathrm{B}(E,F)\) gives the dual \(E\hat{\otimes}F\)
(cf.~\hyperref[section-1-theorem-2]{§1.2, Theorem 2}), we can replace
the weak topology defined by \(E\hat{\otimes}F\) in the ``approximation
problem'' with the bicompact-convergence topology. This also proves that
we can assume \(E\) and \(F\) to be separable. By continuing with such
procedures (notably the systematic use of
\hyperref[section-1-theorem-2]{Theorem 2}), we have given, in {[}PTT,
chap.~1, §5, prop. 37{]}, a large number of other equivalent
formulations of the above conjecture. It suffices, for example, to
assume in the above that \(E\) is a topological-vectorial subspace of
\(c_0\), and that \(F\) is its dual. It even suffices to prove that, if
\(u\in E'\hat{\otimes}E\) defines a zero nuclear operator, then
\(\operatorname{Tr}u=0\). A more concrete formulation of this latter
statement is the following: let \(u=(u_{ij})\) be a matrix that
represents an element of \(\ell^1\hat{\otimes}c_0\) (i.e.~such that
\(\sum_i\sup_j|u_{ij}|<+\infty\)), and such that \(u^2=0\); then
\(\operatorname{Tr}u=\sum u_{ii}=0\). Another formulation: let
\(K(x,y)\) be a continuous kernel on \(X\times X\) (where \(X\) is a
compact space endowed with a positive measure \(\mu\)) such that
\(K\circ K=0\); then
\(\operatorname{Tr}K=\int K(x,x)\mathrm{d}\mu(x)=0\). Or: let \(f(x,y)\)
be a continuous function on the product of two compact spaces \(X\) and
\(Y\); then \(f\) is the uniform limit of linear combinations of
functions of the type \(f(x,b)f(a,y)\). In these latter two examples, it
suffices to do the proof in only \emph{one} case for which the general
conjecture has been proven, provided that \(\mu\) is not the sum of a
sequence of discrete masses in the first case, or that \(X\) and \(Y\)
are infinite in the second case. Other formulations are given in what
follows.

We say that a locally convex space \(E\) satisfies the
\emph{approximation condition} if the identity map from \(E\) to itself
is the limit, in the topology given by uniform convergence on every
precompact subset, of continuous linear maps of finite rank.
\oldpage{96} Then, for \emph{every} locally convex space \(F\),
\emph{every} continuous linear map from \(E\) to \(F\), or from \(F\) to
\(E\), is the limit, in the precompact-convergence topology, of
continuous linear maps of finite rank. If \(E\) is a Banach space, then
this also implies that, for every Banach space \(F\), every compact
linear map from \(F\) to \(E\) is the limit, in the sense of the
\emph{norm}, of continuous linear maps of finite rank; or even that, for
every Banach space \(G\), the space \(G\hat{\hat{\otimes}}E\) is
identical to the space of compact and weakly continuous linear maps from
\(G'\) to \(E\). We can show, by using
\hyperref[section-1-theorem-2]{Theorem 2}, that this is equivalent to
the fact that, for every Banach space \(F\), the canonical map from
\(E\hat{\otimes}F\) to \(\mathrm{B}(E',F')\) is bijective; and, in this
condition, it suffices to take \(F=E'\) (and we thus have here a
\emph{bijectivity condition}), and even to suppose that the trace of any
\(u\in E'\hat{\otimes}E\) that defines a zero operator is zero. --- We
can show that the dual \(E'\) of a Banach space \(E\) satisfies the
approximation condition if and only if every compact linear map from
\(E\) to a Banach space \(F\) is the limit, in the sense of the norm, of
continuous linear maps of finite rank, and that this implies that \(E\)
itself satisfies the approximation condition. --- The ``approximation
problem'' described above can also be stated as follows: does every
locally convex space satisfy the approximation condition? We have
already seen that we can restrict to closed vector subspaces of \(c_0\).

The spaces \(\mathrm{L}^p\) (for \(1\leqslant p\leqslant+\infty\)),
constructed from an arbitrary measure, the spaces \(\mathrm{C}(K)\) (of
continuous functions on a compact space \(K\)), as well as the duals,
biduals, etc. of these spaces, all satisfy the approximation condition
(and even the stronger ``metric approximation condition''; see below).
Nuclear spaces satisfy the approximation condition (see
\hyperref[section-2]{§2}), and so too, more generally, do spaces that
are isomorphic to subspaces of products of Hilbert spaces (these types
of spaces arising rather frequently in practice). I give other examples
in {[}PTT, chap.~1, §5, no. 3{]}, including, notably, the most important
amongst the Banach spaces given by distributions over \(\mathbb{R}^n\)
(essentially those that are between \((\mathscr{D})\) and
\((\mathscr{D}')\) and that are stable under translation and under
multiplication by \(\varphi\in(\mathscr{D})\)).

We say that a Banach space \(E\) satisfies the \emph{metric
approximation condition} if the identity map from \(E\) to itself is the
limit, uniform on every compact subset, of linear applications of finite
rank \emph{and of norm \(\leqslant 1\)}. This is a metric strengthening
of the approximation condition, and we can give analogous reformulations
({[}PTT, chap.~1, §5, no. 2{]}). \oldpage{97} We note the following: for
every Banach space \(F\), the canonical map from \(E\hat{\otimes}F\) to
the space \(\mathrm{J}(E',F')\) of integral bilinear forms on
\(E'\times F'\) is a metric isomorphism. It again suffices to prove that
the canonical map from \(E\hat{\otimes}E'\) to \(\mathrm{J}(E',E)\) is a
metric isomorphism. We do not know of any Banach space that does not
satisfy the metric approximation condition.

We can again prove that \(E'\) satisfies the metric approximation
condition if and only if, for every Banach space \(F\), the canonical
map from \(E'\hat{\otimes}F'\) to the space \(\mathrm{J}(E,F)\) of
integral bilinear forms on \(E\times F\) is a metric isomorphism; we can
then prove that \(E\) then also satisfies the metric approximation
condition. More profound is the following result:

\protect\phantomsection\label{section-1-theorem-9}
\begin{itenv}{Theorem 9}
Let \(E\), \(F\), and \(G\) be Banach spaces, and \(u\) (resp. \(v\)) a
continuous linear map from \(E\) to \(F\) (resp. from \(F\) to \(G\)).
Suppose that one of the maps (\(u\) or \(v\)) is weakly compact, and
that the other is the uniform limit, on every compact subset, of
continuous linear maps of finite rank. Then \(w=v\circ u\) is the
uniform limit, on every compact subset, of continuous linear maps of
finite rank \emph{and of norm \(\leqslant\|w\|\)}.

\end{itenv}

\begin{itenv}{Corollary}
Let \(E\) be a reflexive Banach space. For \(E\) to satisfy the metric
approximation condition, it suffices for it to satisfy the approximation
condition.

\end{itenv}

These statements give conclusions of a metric nature from purely
topological hypotheses, and thus hold true if we replace the given norms
by equivalent norms. In this respect, the corollary to
\hyperref[section-1-theorem-9]{Theorem 9} gives, even for a Hilbert
space, a new approximation result.

\section{Nuclear spaces}\label{section-2}

\subsection{Introduction to nuclear spaces}\label{section-2.1}

\oldpage{98}

In the majority of examples where \(E\) is a given concrete complex
locally convex space (a space of functions, or distributions, for
example), and \(F\) is an arbitrary complete locally convex space, we do
not know how to concretely, simultaneously characterise
\(E\hat{\otimes}F\), \(E\hat{\hat{\otimes}}F\), and its
topological-vectorial extension \(\mathscr{B}_e(E'_s,F'_s)\) (for
example, interpret these spaces as spaces of functions or distributions
with vector values, characterised in a simple way). But we do often know
how to concretely describe a locally convex space \(P\) sitting between
\(E\hat{\otimes}F\) and \(E\hat{\hat{\otimes}}F\), or at least between
\(E\hat{\otimes}F\) and \(\mathscr{B}_e(E'_s,F'_s)\) (``sitting
between'' meaning that the maps \(E\hat{\otimes}F\to P\) and
\(P\to\mathscr{B}_e(E'_s,F'_s)\) are continuous).

\protect\phantomsection\label{section-2-example-1}
\begin{rmenv}{Example 1}
Let \(E=\mathrm{L}^p(\mu)\), where \(1\leqslant p<+\infty\); if \(F\) is
a Banach space (and, by extension, if \(F\) is an aribtrary complex
locally convex space), the we define the space \(\mathrm{L}_F^p(\mu)\)
of ``\(p\)-th power integrable'' maps from the measured space \(M\) to
\(F\) (cf. \autocite{2}) in a natural way, and we easily see that \[
  \mathrm{L}^p\hat{\otimes}F \subset \mathrm{L}_F^p \subset \mathrm{L}^p\hat{\hat{\otimes}}F.
\] In general, \(\mathrm{L}_F^p\) is different from
\(\mathrm{L}^p\hat{\otimes}F\) and \(\mathrm{L}^p\hat{\hat{\otimes}}F\),
and the elements of \(\mathrm{L}^p\hat{\otimes}F\) do not admit an
obvious internal characterisation as maps from \(M\) to \(F\) (unless
\(p=1\), in which case see \hyperref[section-1-theorem-3]{§1, Theorem
3}), whereas the elements of \(\mathrm{L}^p\hat{\hat{\otimes}}F\) can
not even be interpreted, in general, as measurable or scalar-measurable
maps from \(M\) to \(F\).

\end{rmenv}

\protect\phantomsection\label{section-2-example-2}
\begin{rmenv}{Example 2}
For spaces \(E\) consisting of actual functions (and not only of classes
of functions, like \(\mathrm{L}^p\)), we can often characterise
\(\mathscr{B}_e(E'_s,F'_s)\), thanks to the following:

\end{rmenv}

\oldpage{99}

\protect\phantomsection\label{section-1-lemma}
\begin{itenv}{Lemma}
Let \(E\) be a function space on a set \(M\), endowed with a locally
convex topology that is finer than the simple-convergence topology.
Then, for every complete locally convex space \(F\), we can interpret
\(\mathscr{B}_e(E'_s,F'_s)\) as the space of maps \(f\) from \(M\) to
\(F\) such that, for all \(y'\in F'\), the function
\(f_{y'}(t)=\langle f(t),y'\rangle\) belongs to \(E\) (i.e.~\(f\)
belonging \emph{scalar-wise} to \(E\)), and such that \(f_{y'}\) runs
over a weakly relatively compact subset of \(E\) as \(y'\) runs over an
equicontinuous subset of \(F'\). (This second condition is excessive if
\(E\) is reflexive and of type \((\mathscr{F})\) or type
\((\mathscr{LF})\)).

\end{itenv}

But in a case such as the above, we do not, in general, have a handle on
the space \(E\hat{\otimes}F\).

We thus appreciate the simplifications that arise in the case where we
know in advance that \(E\hat{\otimes}F=E\hat{\hat{\otimes}}F\), or even
that \(E\hat{\otimes}F=\mathscr{B}_e(E'_s,F'_s)\) (with these equalities
being assumed to involve the topologies). Then we can often concretely
determine
\(E\hat{\otimes}F=E\hat{\hat{\otimes}}F=\mathscr{B}_e(E'_s,F'_s)\); for
example, if we know a priori an intermediate space \(P\) between
\(E\hat{\otimes}F\) and \(\mathscr{B}_e(E'_s,F'_s)\). At the same time,
we will then have determined the space of weakly continuous linear maps
from \(E'\) to \(F\), or from \(F'\) to \(E\), since these spaces can in
fact be identified with \(\mathscr{B}_e(E'_s,F'_s)\). And the dual
\(\mathrm{B}(E,F)\) of \(E\hat{\otimes}F\), i.e.~the space of continuous
bilinear forms on \(E\times F\), will then be identical to the dual of
\(E\hat{\hat{\otimes}}F\), or also to the dual of \(P\), and can thus
often be concretely understood.

For example, the essence of L. Schwartz's ``kernel theorem'' says that
the space of continuous bilinear forms on
\(\mathscr{E}(\mathbb{R}^m)\times\mathscr{E}(\mathbb{R}^m)\) is
identical to the space of bilinear forms defined by the distributions on
\(\mathbb{R}^m\times\mathbb{R}^n\) with compact support. The latter are
a priori the continuous linear forms on
\(\mathscr{E}(\mathbb{R}^m\times\mathbb{R}^n)\), but we directly see (as
a particular case of the above lemma) that
\(\mathscr{E}(\mathbb{R}^m\times\mathbb{R}^n)\) can be identified with
the space \[
  E\hat{\hat{\otimes}}F = \mathscr{B}_e(E'_s,F'_s)
\] where \(E=\mathscr{E}(\mathbb{R}^m)\) and
\(F=\mathscr{E}(\mathbb{R}^n)\), and so a priori the distribution on
\(\mathbb{R}^m\times\mathbb{R}^n\) are the \emph{integral}
(\hyperref[section-1.8]{§1.8}) bilinear forms on
\(\mathscr{E}(\mathbb{R}^m)\times\mathscr{E}(\mathbb{R}^m)\). The kernel
theorem, which says that \emph{all} continuous bilinear forms on
\(\mathscr{E}(\mathbb{R}^m)\times\mathscr{E}(\mathbb{R}^m)\) are
obtained in this way, i.e.~that the transpose of the canonical map from
\(\mathscr{E}(\mathbb{R}^m)\hat{\otimes}\mathscr{E}(\mathbb{R}^m)\) to
\(\mathscr{E}(\mathbb{R}^m)\hat{\hat{\otimes}}\mathscr{E}(\mathbb{R}^m)\)
is a \emph{surjective} map, also implies (by a classical theorem of
\((\mathscr{F})\)-spaces) that we in fact have that \[
  \mathscr{E}(\mathbb{R}^m)\hat{\otimes}\mathscr{E}(\mathbb{R}^m)
  = \mathscr{E}(\mathbb{R}^m)\hat{\hat{\otimes}}\mathscr{E}(\mathbb{R}^m).
\] \oldpage{100} In fact, the proof of the theorem even shows that \[
  \mathscr{E}(\mathbb{R}^m)\hat{\otimes}F
  = \mathscr{E}(\mathbb{R}^m)\hat{\hat{\otimes}}F
\] (a space which is isomorphic, when \(F\) is complete, to the space of
infinitely differentiable maps from \(\mathbb{R}^m\) to \(F\)) for
\emph{every} locally convex space \(F\).

We thus see that, in general, \emph{the claim that
\(E\hat{\otimes}F=E\hat{\hat{\otimes}}F\) for two given spaces \(E\) and
\(F\) should be seen as an algebraic-topological equivalent of L.
Schwartz's kernel theorem}. In what follows, we will study the spaces
\(E\), such as \(\mathscr{E}(\mathbb{R}^m)\), for which we have
\(E\hat{\otimes}F=E\hat{\hat{\otimes}}F\) for \emph{every} locally
convex space \(F\): these are the \emph{nuclear spaces}.

\subsection{Characterisation of nuclear spaces}\label{section-2.2}

{[}PTT, chap.~2, §2, no. 1{]}

By \hyperref[section-2.1]{§2.1}, a locally convex space is said to be
\emph{nuclear} if, for every locally convex space \(F\), the canonical
map from \(E\hat{\otimes}F\) to \(E\hat{\hat{\otimes}}F\) is a
topological-vectorial isomorphism from the first space to the second
(or, equivalently, if these two spaces induce the same topology on
\(E\otimes F\)). It actually suffices to check this condition only in
the case where \(F\) is a Banach space. Further, \(E\) is nuclear if and
only if its completion is nuclear.

\protect\phantomsection\label{section-2-theorem-1}
\begin{itenv}{Theorem 1}
Let \(E\) be a locally convex space. Then the following conditions are
equivalent:

\begin{enumerate}
\def\labelenumi{\alph{enumi}.}
\item
  \(E\) is nuclear.
\item
  Every continuous linear map from \(E\) to a Banach space \(F\) is
  nuclear ((§1.7){[}\#section-1.7{]}).
\item
  For every equicontinuous weakly closed disk \(A\) in \(E'\), there
  exists another \(B\supset A\) such that the identity map from \(E'_A\)
  to \(E'_B\) is nuclear.
\end{enumerate}

If \(E\) is then complete, we have, for every complete locally convex
space \(F\), that \[
  E\hat{\otimes}F = E\hat{\hat{\otimes}}F = \mathscr{B}_e(E'_s,F'_s)
\] (i.e.~topological-vectorial isomorphisms).

\end{itenv}

\(\star\) The hard part of the proof consists of showing that (a)
implies (c). By taking the dual of condition (a), we find that, for
every equicontinuous weakly closed disk \(A\) of \(E'\), there exists
another \(B\supset A\) such that the identity map from \(E'_A\) to
\(E'_B\) is \emph{integral} ((§1.8){[}\#section-1.8{]}). Applying the
same result to \(B\), and using the corollary of
\hyperref[section-1-theorem-8]{Theorem 8}, we find (c). We note that we
have only used the fact that, for every \emph{Banach} space \(F\), the
canonical map from \(E\hat{\otimes}F\) to \(E\hat{\hat{\otimes}}F\) is a
\emph{weak} isomorphism. \(\star\)

\oldpage{101}

Since a nuclear map is a fortiori compact, (c) implies the following:

\protect\phantomsection\label{section-2-corollary-1}
\begin{itenv}{Corollary 1}
Let \(E\) be a nuclear space. Then the bounded subsets of \(E\) are
precompact (and \(E\) is even a Schwartz space, cf. \autocite{6}). A
fortiori, if \(E\) is quasi-complete, then \(E\) is of type
\((\mathscr{M})\), and thus reflexive.

\end{itenv}

We thus also conclude that a nuclear Banach space is of finite
dimension. Furthermore:

\protect\phantomsection\label{section-2-corollary-2}
\begin{itenv}{Corollary 2}
Let \(E\) be a nuclear space, and \(F\) an arbitrary locally convex
space. If \(F\) is quasi-complete, then every bounded linear map from
\(E\) to \(F\) is nuclear. If \(E\) is barrelled (e.g.~if \(E\) is of
type \((\mathscr{F})\)), then every bounded linear map from \(F\) to
\(E'\) is nuclear.

\end{itenv}

\protect\phantomsection\label{section-2-corollary-3}
\begin{itenv}{Corollary 3}
Let \(E\) and \(F\) be locally convex spaces, with \(E\) nuclear. Then
the continuous bilinear forms on \(E\times F\) are exactly the
``nuclear'' bilinear forms, i.e.~those which come from an element of a
space of the form \(E'_A\hat{\otimes}F'_B\), where \(A\) (resp. \(B\))
is an equicontinuous weakly closed disk in \(E'\) (resp. \(F'\)).

\end{itenv}

We do not know of any case where we have
\(E\hat{\otimes}F=E\hat{\hat{\otimes}}F\) without having either \(E\) or
\(F\) already nuclear (in this problem, we can easily reduce to the case
where \(E\) and \(F\) are both of type \((\mathscr{F})\)). We can show
that, if \(F\) is the space \(c_0\) or \(\ell^1\), or, more generally,
if \(F\) admits a quotient space isomorphic to either \(c_0\) or
\(\ell^1\), then the statement that
\(E\hat{\otimes}F=E\hat{\hat{\otimes}}F\) already implies that \(E\) is
nuclear. Combining this with Banach's classical ``isomorphism theorem'',
we obtain:

\protect\phantomsection\label{section-2-theorem-2}
\begin{itenv}{Theorem 2}
Let \(E\) be an \((\mathscr{F})\)-space. For \(E\) to be nuclear, it is
necessary and sufficient for every summable sequence in \(E\) to be
absolutely summable, or also for every sequence in \(E\) that converges
to \(0\) to converge to \(0\) nuclearly (cf.~the end of
(§2.4){[}\#section-2.4{]}).

\end{itenv}

In particular, if \(E\) is a Banach space, then each of the two above
conditions imply that \(E\) is of finite dimension: in the first case,
we thus recover a recent theorem of Dvoretzky--Rogers.

\subsection{Invariance properties, examples of nuclear
spaces}\label{section-2.3}

{[}PTT, chap.~2, §2, no. 2 and no. 3{]}

The class of nuclear spaces enjoys a remarkable stability property,
expressed in the following:

\oldpage{102}

\protect\phantomsection\label{section-2-theorem-3}
\begin{itenv}{Theorem 3}

---

\begin{enumerate}
\def\labelenumi{\arabic{enumi}.}
\tightlist
\item
  Let \(E\) be an \((\mathscr{F})\)-space or a \((\mathscr{DF})\)-space.
  Then, for \(E\) to be nuclear, it is necessary and sufficient for
  \(E'\) to be nuclear.
\item
  Let \(E\) be a nuclear space. Then every vector subspace and every
  quotient space of \(E\) is nuclear.
\item
  The topological-vectorial product of an arbitrary family of nuclear
  spaces, and the topological-vectorial sum of a countable family of
  nuclear spaces, is nuclear.
\item
  Let \(E\) and \(F\) be two nuclear spaces. Then \(E\hat{\otimes}F\) is
  nuclear.
\item
  Let \(E\) be a nuclear space of type \((\mathscr{F})\), and \(F\) a
  reflexive space whose dual is nuclear. Then the dual of
  \(E\hat{\otimes}F\) is nuclear.
\end{enumerate}

\end{itenv}

\(\star\) Of these claims, (2) and (5) are the most difficult (the
others follow reasonably simply from the criteria of
\hyperref[section-2-theorem-1]{Theorem 1}). The proof given in
{[}PTT{]}, which relies on rather acute techniques, can be replaced by
the systematic use of criterion (c) of
\hyperref[section-2-theorem-1]{Theorem 1}, along with the two following
facts:

\begin{enumerate}
\def\labelenumi{\arabic{enumi}.}
\tightlist
\item
  If \(E\) is a nuclear space, then there exists a fundamental family of
  equicontinuous subsets \(A\) of \(E'\) such that \(E'_A\) is a Hilbert
  space (i.e.~\(E\) is isomorphic to a vector subspace of the product of
  a family of Hilbert spaces);
\item
  The remark made at the end of \hyperref[section-1.7]{§1.7}. \(\star\)
\end{enumerate}

From \hyperref[section-2-theorem-3]{Theorem 3}, we obtain the following
results: if \(E\) is a vector space endowed with the finest topology
making some family of linear maps \(f_i\) from \(E\) to nuclear spaces
\(E_i\) continuous, then \(E\) is nuclear; if \(E\) is a vector space,
and \(u_i\) are linear maps from a sequence of nuclear spaces \(E_i\) to
\(E\) such that \(E\) is generated by the union of the \(u_i(E_i)\), and
if endow \(E\) with the finest locally convex topology making the
\(u_i\) continuous, then \(E\) is nuclear; if \(E\) is a nuclear space
of type \((\mathscr{F})\) or \((\mathscr{DF})\), and \(F\) is a nuclear
space, then \(\mathrm{L}_b(E,F)\) is nuclear, and so too is the dual of
\(\mathrm{L}_b(E,F)\) if \(F\) is also of type \((\mathscr{F})\) or
\((\mathscr{DF})\).

\hyperref[section-2-theorem-3]{Theorem 3}, along with the fact that the
space \(\mathscr{E}(U)\) of infinitely differentiable functions on an
open subsets \(U\) of \(\mathbb{R}^n\) is nuclear
(\hyperref[section-2.1]{§2.1}), allows us to easily prove the following:

\begin{itenv}{Corollary}
All of the following are nuclear: the spaces \(\mathscr{E}\),
\(\mathscr{E}'\), \(\mathscr{D}\), and \(\mathscr{D}'\), constructed on
an open subset \(U\) of \(\mathbb{R}^n\); the spaces \((\mathscr{S})\),
\((\mathscr{S}')\), \((\mathscr{O}_M)\), and \((\mathscr{O}'_C)\),
constructed on \(\mathbb{R}^n\); and the strong duals of these latter
two spaces. The space \(\mathscr{H}\) of holomorphic functions on a
complex-analytic manifold is also nuclear. (For the definitions of the
above spaces, see \autocite{9}).

\end{itenv}

\oldpage{103}

\(\star\) Consider an increasing sequence of positive sequences
\(a^{(n)}=(a_i^{(n)})\), and let \(E\) be the corresponding ``echelon''
space (see \autocite{7}), i.e.~the space of sequences \((x_i)\) such
that \(\sum_i x_ia_i^{(n)}\) is absolutely summable for all \(n\). Note
that \(E\) is endowed with a natural topology of an \((\mathscr{F})\)
space. \emph{For \(E\) to be nuclear, it is necessary and sufficient for
there to exist, for all \(n\), some \(m\geqslant n\) such that the
sequence \(a^{(n)}/a^{(m)}=(a_i^{(n)}/a_i^{(m)})\) to be summable}
(where we take a quotient of two zero terms to be zero). Thus the space
\((s)\) of rapidly decreasing sequences, as well as its dual \((s')\)
(the space of slowly increasing sequences), is nuclear. (This easily
recovers, thanks to the Fourier transform, the fact that
\(\mathscr{E}(\mathbb{R}^n)\), and, more generally, \(\mathscr{E}(U)\),
is nuclear.) The above statement also allows us to construct echelon
spaces that are of type \((\mathscr{M})\) (and even of type
\((\mathscr{S})\), see \autocite[Section 3]{6}) but not nuclear.
\(\star\)

\subsection{Lifting properties}\label{section-2.4}

{[}PTT, chap.~2, §3, no. 1{]}

\hyperref[section-1-theorem-6]{Theorem 6 of §1.7} can be specialised as
follows:

\protect\phantomsection\label{section-2-theorem-4}
\begin{itenv}{Theorem 4}

---

\begin{enumerate}
\def\labelenumi{\arabic{enumi}.}
\item
  Let \(E_1\) and \(E_2\) be locally convex spaces, and \(F_1\) (resp.
  \(F_2\)) a vector subspace of \(E_1\) (resp. \(E_2\)). Suppose that
  \(F_1\) or \(F_2\) is nuclear. Then every continuous bilinear form on
  \(F_1\times F_2\) is the restriction of a nuclear bilinear form on
  \(E_1\times E_2\).
\item
  Let \(E\) and \(G\) be locally convex spaces, and \(F\) a vector
  subspace of \(E\). Suppose that \(F\) is nuclear and that \(G\) is
  quasi-complete, or that \(G\) is the strong dual of a quasi-barrelled
  nuclear space (for example, that \(G\) is a complete nuclear space of
  type \((\mathscr{F})\) or \((\mathscr{DF})\)). Then every bounded
  linear map from \(F\) to \(G\) is the restriction of a nuclear map
  from \(E\) to \(G\).
\item
  Let \(E\) and \(G\) be locally convex spaces, and \(F\) a closed
  vector subspace of \(E\), such that every compact disk of \(E/F\) is
  contained in the canonical image of a bounded complete disk of \(E\).
  Suppose that \(G\) is nuclear, or that \(E/F\) is isomorphic to the
  strong dual of a quasi-barrelled nuclear space (for example, that
  \(E\) is a complete nuclear space of type \((\mathscr{F})\) or
  \((\mathscr{DF})\)). Then every bounded linear map from \(G\) to
  \(E/F\) comes from a nuclear map from \(G\) to \(E\).
\end{enumerate}

\end{itenv}

A variant of the above properties is the following: Let \(E\) be a
locally convex space, \(F\) a nuclear vector subspace, and \(A\) an
equicontinuous disk in \(E'\); then there exists a linear map from
\(F'_A\) to \(E'\) that is right inverse to the canonical map from
\(E'\) to \(F'\), and that sends \(A\) to an equicontinuous subset of
\(E'\). \oldpage{104} We note again the following ``lifting'' property,
which is an easy consequence of the above results:

\protect\phantomsection\label{section-2-theorem-5}
\begin{itenv}{Theorem 5}
Let \(E\) and \(F\) be locally convex spaces, both of type
\((\mathscr{F})\), or both of type \((\mathscr{DF})\). Suppose that
\(E\) is nuclear, and let \(A\) be a bounded disk in
\(P=E\hat{\otimes}F\). Then there exists a sequence \((x_i)\) that tends
to \(0\) in \(E\), an equicontinuous sequence of linear maps
\(u\mapsto y_i(u)\) from \(P_A\) to \(F\), and a fixed summable sequence
\((\lambda_i)\), such that \[
  u = \sum \lambda_i x_i\otimes y_i(u)
\] for all \(u\in P_A\).

\end{itenv}

\begin{itenv}{Corollary}
The bounded disk \(A\) is contained in the closed disked hull of a set
\(B\otimes C\), where \(B\) (resp. \(C\)) is some bounded subset of
\(E\) (resp. \(F\)). Thus, on \(\mathrm{B}(E,F)\), the bibounded
convergence topology and the topology of the strong dual of
\(E\hat{\otimes}F\) are identical.

\end{itenv}

We note again that the analogue of
\hyperref[section-2-theorem-5]{Theorem 5} for the representation of
\emph{convergent sequences} in \(E\hat{\otimes}F\) also holds true.

\subsection{The tensor product of a nuclear space with an arbitrary
locally convex space}\label{section-2.5}

{[}PTT, chap.~2, §3, no. 2 and no. 3{]}

The duality theory for \(E\hat{\otimes}F\) is particularly simple when
\(E\) is nuclear, and when \(E\) and \(F\) are both of type
\((\mathscr{F})\), or both of type \((\mathscr{DF})\). In fact, the
above results easily give:

\protect\phantomsection\label{section-2-theorem-6}
\begin{itenv}{Theorem 6}
Let \(E\) and \(F\) be locally convex spaces, both of type
\((\mathscr{F})\), or both of type \((\mathscr{DF})\). If \(E\) is
nuclear, then the strong dual of \(E\hat{\otimes}F\), which can be
identified with \(\mathrm{B}(E,F)\) endowed with the bibounded
convergence topology (cf.~the corollary of
\hyperref[section-2-theorem-5]{Theorem 5}), can also be identified with
\(E'\hat{\otimes}F'\).

\end{itenv}

Subsequently, if \(E\) is complete, and thus reflexive, then the bidual
of \(E\hat{\otimes}F\), endowed with the strong topology of the dual of
\((E\otimes F)'\), can be identified with \(E\hat{\otimes}F''_b\), where
\(F''_b\) denotes the strong dual of \(F'\).

In general, if \(E\) is nuclear and complete, but \(E\) and \(F\) are
otherwise arbitrary, and if we again take the ``natural'' topologies on
the biduals (cf.~Introduction), then we can prove that the bidual of
\(E\hat{\otimes}F\) can be identified with a dense topological-vectorial
subspace of \(E\hat{\otimes}F''\) (see {[}PTT, chap.~1, §4, no. 2, prop.
25{]}, where we give a more general statement). As for the dual of
\(E\hat{\otimes}F\), whose elements are characterised by
\hyperref[section-2-corollary-3]{Corollary 3 of Theorem 1}, it will, in
general, no longer be identical to \(E'\hat{\otimes}F'\) at all, but,
under rather strong conditions (such as if \(F\) is reflexive), it will
be a dense topological-vectorial subspace of \(E'\bar{\otimes}F\) (see
the definition in \hyperref[section-1-appendix-1]{§1, Appendix 1}).
\oldpage{105} But even in the case of the tensor product
\(E'\hat{\otimes}E\), where \(E\) is nuclear and of type
\((\mathscr{F})\) or type \((\mathscr{DF})\), the strong dual of
\(E\hat{\otimes}F\) will often often not be complete (and thus distinct
from \(E'\bar{\otimes}F'\)): see \hyperref[section-2-appendix]{the
appendix}.

The product \(E\hat{\otimes}F\) (with \(E\) nuclear) gives rise to many
invariance theorems: for \(E\hat{\otimes}F\) to be reflexive (resp. of
type \((\mathscr{M})\), resp. of type \((\mathscr{S})\), resp.
quasi-normable), it suffices (and is necessary, if \(F\) is complete)
for \(F\) to be so. But it can be the case that \(E\) and \(F\) are
bornological and barrelled without \(E\hat{\otimes}F\) being so (see
\hyperref[section-2-appendix]{the appendix}).

By the comments of \hyperref[section-2.1]{§2.1}, concretely determining
\(E\hat{\otimes}F\) when \(E\) is nuclear usually poses no difficulty.
We now give some examples (where \(F\) is assumed to be a complete
locally convex space).

Let \(\mathscr{E}=\mathscr{E}(V)\) be the space of infinitely
differentiable functions on some given infinitely differentiable
manifold \(V\). Then \(\mathscr{E}\hat{\otimes}F\) is the space of
infinitely differentiable functions from \(V\) to \(F\) (endowed with
its natural topology). The analogous statement holds when we replace
\(\mathscr{E}\) by the subspace \(E\) of functions satisfying some given
system of partial differential equations \(\mathrm{D}_if=0\).

Let \(\mathscr{H}=\mathscr{H}(V)\) be the space of holomorphic functions
on some given complex-analytic manifold \(V\). Then
\(\mathscr{H}\hat{\otimes}F\) is the space of holomorphic maps from
\(V\) to \(F\) (endowed with the compact convergence topology).

In particular, we find that \[
  \begin{aligned}
    \mathscr{E}(U)\hat{\otimes}\mathscr{E}(V) &= \mathscr{E}(U\times V)
  \\\mathscr{H}(U)\hat{\otimes}\mathscr{H}(V) &= \mathscr{H}(U\times V)
  \end{aligned}
\] where \(U\) and \(V\) are some given infinitely differentiable (resp.
holomorphic) manifolds.

Let \(\mathscr{S}=\mathscr{S}(\mathbb{R}^n)\) be the space of
``infinitely differentiable and rapidly decreasing'' functions on
\(\mathbb{R}^n\) (cf. \autocite[t. 2]{9}). Then
\(\mathscr{S}\hat{\otimes}F\) is the space of maps from \(\mathbb{R}^n\)
to \(F\) that are ``infinitely differentiable and rapidly decreasing'',
i.e.~infinitely differentiable, and such that, for every derivation
multi-index \(\mathrm{D}^p\), and for every polynomial \(P\) on
\(\mathbb{R}^n\), the function \(P\mathrm{D}^pf\) is bounded on
\(\mathbb{R}^n\). In particular, \[
  \mathscr{S}(\mathbb{R}^m)\hat{\otimes}\mathscr{S}(\mathbb{R}^n) = \mathscr{S}(\mathbb{R}^m\times\mathbb{R}^n).
\]

We have analogous statements given by taking \(E=(\mathscr{O}_M)\) (the
space of ``infinitely differentiable and slowly increasing'' functions)
or \(E=(\mathscr{O}_C)\) (the space of ``infinitely differentiable and
very slowly increasing'' functions, the strong dual of the space
\((\mathscr{O}'_C)\) of rapidly decreasing distributions \autocite{9}),
etc. \oldpage{106} We also note that, by using the
\hyperref[section-1-lemma]{Lemma in §1}, we find that a map \(f\) with
values in \(F\) is infinitely differentiable (resp. holomorphic, resp.
``infinitely differentiable and rapidly decreasing'', resp. ``infinitely
differentiable and slowly increasing'', etc.) if and only if it is so
``scalar-wise'', i.e.~if and only if, for all \(y'\in F'\), the
scalar-valued function \(\langle f(t),y'\rangle\) has the property in
question.

\subsection{\texorpdfstring{\(p\)-th power summable operators,
application to nuclear
spaces}{p-th power summable operators, application to nuclear spaces}}\label{section-2.6}

{[}PTT, chap.~2, §1, no. 1 to no. 6, and §2, no. 4{]}

Let \(E\) and \(F\) be locally convex spaces. Among the elements of of
\(E\bar{\otimes}F\) (see the definition in
\hyperref[section-1-appendix-1]{§1, Appendix 1}), we have given the name
of \emph{Fredholm kernels} to those which come from an element of some
space \(E_A\hat{\otimes}F_B\), where \(A\) (resp. \(B\)) is a bounded
disk in \(E\) (resp. \(F\)) such that \(E_A\) (resp. \(F_B\)) is
complete; we can then even suppose \(A\) and \(B\) to be compact
(\hyperref[section-1-appendix-1]{§1, Appendix 1}). More generally, if
\(0<p\leqslant 1\), then we call any element of \(E\bar{\otimes}F\) of
the form \(\sum\lambda_i x_i\otimes y_i\), where
\((\lambda_i)\in\ell^p\) and \((x_i)\) (resp. \((y_i)\)) is a sequence
inside a compact subset of \(E\) (resp. \(F\)), a \emph{\(p\)-th power
summable Fredholm kernel in \(E\bar{\otimes}F\)}. (If \(p=1\), then we
recover Fredholm kernels). The set of \(p\)-th power summable Fredholm
kernels in \(E\bar{\otimes}F\) is denoted by
\(E\overset{(p)}{\otimes}F\).

We define a \emph{\(p\)-th power summable map} from \(E\) to \(F\) to be
any linear map from \(E\) to \(F\) that comes from a \(p\)-th power
summable Fredholm kernel in \(E'\bar{\otimes}F\) (if \(p=1\), then we
recover Fredholm maps, introduced in \hyperref[section-1-appendix-1]{§1,
Appendix 1}). The space of these maps, denoted by
\(\mathrm{L}^{(p)}(E,F)\), can thus be identified with a quotient space
of \(E'\overset{(p)}{\otimes}F\), and even with
\(E'\overset{(p)}{\otimes}F\) in all known cases (see the ``problem of
bijectivity'' in \hyperref[section-1-appendix-2]{§1, Appendix 2}); this
is always the case, for example, if \(p\leqslant 2/3\). When \(E\) and
\(F\) are Banach spaces, we can introduce a ``distance to the origin''
function on \(E\otimes F\), by \[
  S_p(u) = \inf\sum|\lambda_i|^p
\] where the infimum is taken over all representations
\(u=\sum\lambda_i x_i\otimes y_i\) with \(\|x_i\|,\|y_i\|\leqslant 1\).
Then \(E\overset{(p)}{\otimes}F\) becomes a complete metrisable
topological-vectorial space which is, in general, \emph{not locally
convex}. When \(E\) and \(F\) are arbitrary locally convex spaces, the
semi-norms of \(E\) and \(F\) allow us to define a system of ``distances
to the origin'' on \(E\overset{(p)}{\otimes}F\), which again make a
topological-vectorial space (which is, in general, not locally convex);
\oldpage{107} this space is metrisable and complete if \(E\) and \(F\)
are metrisable and complete. Consequently, if \(E\) is of type
\((\mathscr{DF})\), and \(F\) of type \((\mathscr{F})\), then
\(\mathrm{L}^{(p)}(E,F)\) is a complete metrisable topological-vectorial
space (which is, in general, not locally convex) isomorphic to a
quotient space of \(E'\overset{(p)}{\otimes}F\).

We can also, for \(0\leqslant p<1\), introduce the space
\(E\overset{[p]}{\otimes}F\) of \emph{Fredholm kernels of order
\(\leqslant p\)}, which is defined to be the intersection of the spaces
\(E\overset{(q)}{\otimes} F\) for \(p<q\leqslant 1\). Endowed with the
topology that is the upper bound of the topologies induced by the spaces
\(E\overset{(q)}{\otimes} F\), this is a topological-vectorial space
(which is not, in general, locally convex), which is metrisable and
complete if \(E\) and \(F\) are metrisable and complete. We analogously
define the topological-vectorial spaces \(\mathrm{L}^{[p]}(E,F)\) as the
intersection of the spaces \(\mathrm{L}^{(q)}(E,F)\) for
\(p<q\leqslant 1\). The introduction of these topologies allows us to
prove the following: let \(E\) and \(F\) both be
\((\mathscr{F})\)-spaces; then every element of
\(E\overset{[p]}{\otimes}F\) (where \(0\leqslant p<1\)) is of the form
\(\sum\lambda_i x_i\otimes y_i\), where \((\lambda_i)\in\ell^{[p]}\),
and where \((x_i)\) (resp. \((y_i)\)) is a bounded sequence in \(E\)
(resp. \(F\)); further, if \((u_\alpha)\) is a sequence that tends to
\(0\) in \(E\overset{(p)}{\otimes}F\) (resp. in
\(E\overset{[p]}{\otimes}F\)), then
\(u_\alpha=\sum\lambda_i^\alpha x_i\otimes y_i\), where the bounded
sequences \((x_i)\) and \((y_i)\) are fixed, and where
\(\lambda^\alpha=(\lambda_i^\alpha)\) tends to \(0\) in \(\ell^p\)
(resp. \(\ell^{[p]}\)). (We denote by \(\ell^{[p]}\) the space given by
the intersection of the spaces \(\ell^q\) for \(p<q\leqslant 1\),
endowed with the topology that is the upper bound of the topologies
induced by the \(\ell^q\), which in fact gives a complete metrisable
space). The most important application of these results is the
following:

\protect\phantomsection\label{section-2-theorem-7}
\begin{itenv}{Theorem 7}
Let \(E\) be a \((\mathscr{DF})\)-space, and \(F\) an
\((\mathscr{F})\)-space. Then every map of order \(0\) from \(E\) to
\(F\) is of the form \(\sum\lambda_i x'_i\otimes y_i\), where
\((\lambda_i)\) is a \emph{rapidly decreasing} sequence of scalars, and
where \((x'_i)\) (resp. \((y_i)\)) is a sequence that tends to \(0\) in
\(E'\) (resp. \(F\)).

\end{itenv}

Let \(E\) be an arbitrary locally convex space. Then every Fredholm
kernel element of \(E'\bar{\otimes}E\) admits a \emph{Fredholm
determinant} \(\det(1-zu)\), which is an entire function in the complex
variable \(z\), and whose zeros are the inverses of the eigenvalues
\(z_i\) of the Fredholm operator defined by \(u\) (where we count the
zeros and the eigenvalues with multiplicity). We have the following:

\oldpage{108}

\protect\phantomsection\label{section-2-theorem-8}
\begin{itenv}{Theorem 8}
If \(u\in E'\overset{(p)}{\otimes}E\), then the Fredholm determinant of
\(u\) is an entire function of order \(\leqslant q\), where
\(1/q=1/p-1/2\), and the sequence of eigenvalues of \(u\) is \(p\)-th
power summable. If \(p\leqslant 2/3\) (and so \(q\leqslant 1\)), then
\(\det(1-zu)\) is an entire function \emph{of genus \(0\)}, and thus
identical to the infinite product \(\prod(1-zz_i)\) (where the \(z_i\)
are the eigenvalues of \(u\), repeated according to their order of
multiplicity). If \(E\) is a Banach space, then we have \[
  \left(\sum|z_i|^q\right)^{\frac1q} \leqslant(S_p(u))^{\frac1p}.
\]

\end{itenv}

\protect\phantomsection\label{section-2-theorem-8-corollary-1}
\begin{itenv}{Corollary 1}
The sequence of eigenvalues of a Fredholm operator is square summable.

\end{itenv}

\protect\phantomsection\label{section-2-theorem-8-corollary-2}
\begin{itenv}{Corollary 2}
If \(u\) is a \(\frac23\)-summable Fredholm operator, then
\(\operatorname{Tr}u=\sum_i z_i\) (where the \(z_i\) are the eigenvalues
of \(u\)).

\end{itenv}

\protect\phantomsection\label{section-2-theorem-8-corollary-3}
\begin{itenv}{Corollary 3}
The sequence of eigenvalues of a Fredholm operator of order \(0\) in
\(E\), arranged in order of decreasing modulus, is a rapidly decreasing
sequence.

\end{itenv}

Furthermore, we can show that, under the conditions of
\hyperref[section-2-theorem-8]{Theorem 8}, when \(E\) is a Banach space,
the sequence of eigenvalues of \(u\in E'\overset{(p)}{\otimes}E\), as an
unordered \(q\)-th power summable sequence, depends continuously on
\(u\). This implies that, if \(u_\alpha\to u\) in
\(E'\overset{(p)}{\otimes}E\), then we can arrange the eigenvalues of
\(u_\alpha\) into a sequence
\(\lambda^{(\alpha)}=(\lambda_i^{(\alpha)})\) such that
\(\lambda^{(\alpha)}\to\lambda^{(0)}\) in \(\ell^q\). We have an
analogous result for \(E'\overset{[p]}{\otimes}E\).

By iterating Fredholm kernels, we obtain kernels that have stronger and
stronger ``decreasing properties'':

\protect\phantomsection\label{section-2-theorem-9}
\begin{itenv}{Theorem 9}
Let \(u\) be a Fredholm kernel given by the composition of \(n\)
Fredholm kernels \(u_i\), with each \(u_i\) being \(p_i\)-th power
summable (with \(0<p_i\leqslant 1\)). Then \(u\) is \(p\)-th power
summable, where \[
  \frac1p = \left(\sum\frac{1}{p_i}\right)-\frac{n+1}{2}.
\] If \(u_i\in E'_i\overset{(p_i)}{\otimes}E_{i+1}\), where the \(E_i\)
are Banach spaces, then \[
  S_p(u)^{\frac1p} \leqslant\prod\left(S_{p_i}(u_i)\right)^{\frac{1}{p_i}}.
\]

\end{itenv}

In any case, the composition of \(n\) Fredholm operators (for
\(n\geqslant 3\)) is is \((2/(n-1))\)-th power summable. For \(n=2\) or
\(n=3\), the formula \(1/p=(\sum(1/p_i))-(n+1)/2\) does not give any
information, but it seems that we should replace \(n+1\) by \(n-1\) in
it.

\oldpage{109}

For more details on the above classes of operators, I refer the reader
to {[}PTT, chap.~2, §1{]}. We note only that we can show that, up to a
small numerical uncertainty in the value of the exponents (uncertainty
being the very nature of things), the fact that, for an operator to be
\(p\)-th power summable, it suffices to study the properties of the
image of a suitable neighbourhood of \(0\) in \(E\) (the unit ball, if
\(E\) is a Banach space). The uncertainty disappears for operators of
order \(0\), which are characterised exactly by the fact that they send
a suitable neighbourhood of \(0\) in \(E\) to a subset of \(F\) that is
``of order \(0\)'' {[}PTT, chap.~2, §1, no. 5{]}. When \(F\) is of type
\((\mathscr{F})\), the subsets of \(F\) of order \(0\) are the subsets
contained in the closed disked hull of a ``rapidly decreasing'' sequence
in \(F\).

The combination of \hyperref[section-2-theorem-1]{Theorem 1} and
\hyperref[section-2-theorem-9]{Theorem 9} gives:

\protect\phantomsection\label{section-2-theorem-10}
\begin{itenv}{Theorem 10}
Let \(E\) be a nuclear space. Then, for every \(\varepsilon>0\), and
every weakly closed equicontinuous disk \(A\) in \(E'\), there exists a
weakly closed equicontinuous disk \(B\supset A\) such that the identity
map from \(E'_A\) to \(E'_B\) is \(\varepsilon\)-th power summable.

\end{itenv}

\protect\phantomsection\label{section-2-theorem-10-corollary-1}
\begin{itenv}{Corollary 1}
Let \(E\) be a nuclear space, and \(F\) a locally convex space. Then
every quasi-complete bounded linear map from \(E\) to \(F\), and every
linear map from \(F\) to \(E'\) that sends a suitable neighbourhood of
\(0\) to an equicontinuous subset, is of order \(0\). Every continuous
bilinear form on \(E\times F\) comes from an element of
\(E'\overset{[0]}{\otimes}F'\) (i.e.~is a Fredholm kernel of order
\(0\)).

\end{itenv}

\protect\phantomsection\label{section-2-theorem-10-corollary-2}
\begin{itenv}{Corollary 2}
Let \(E\) be a quasi-complete nuclear space. Then every bounded operator
to \(E\) is a Fredholm operator of order \(0\), whose Fredholm
determinant is thus of order \(0\), and the sequence of eigenvalues,
arranged in order of decreasing modulus, is a rapidly decreasing
sequence.

\end{itenv}

Other corollaries can be obtained by taking into account the other
results given in this section. Thus, if \(E\) and \(F\) are both of type
\((\mathscr{F})\), with \(E\) nuclear, then every element of
\(E\hat{\otimes}F\) is of the form \(\sum\lambda_i x_i\otimes y_i\),
where \((x_i)\) (resp. \((y_i)\)) is a bounded sequence in \(E\) (resp.
\(F\)), and where \((\lambda_i)\) is a \emph{rapidly decreasing}
sequence. If \(E\) and \(F\) are of type \((\mathscr{DF})\), with \(E\)
still nuclear, then we only assume that
\(\sum|\lambda_i|^\varepsilon<+\infty\), with \(\varepsilon>0\) being
given in advance. Similarly, \hyperref[section-2-theorem-5]{Theorem 5}
can be made more precise in an analogous way. More generally, if \(E\)
and \(F\) are both locally convex spaces, with \(E\) nuclear, then every
continuous bilinear form \(u\) on \(E\times F\) is of the form
\(\sum\lambda_i x'_i\otimes y'_i\), where
\((\lambda_i)\in\ell^\varepsilon\), and where \((x'_i)\) (resp.
\((y'_i)\)) is an equicontinuous sequence in \(E'\) (resp. \(F'\)).
\oldpage{110} If \(u\) varies in a weakly closed equicontinuous disk
\(M\) of bilinear forms, then we can take \((x'_i)\) and \((\lambda_i)\)
in the above to be fixed, and \(y'_i=\varphi_i(u)\), where the
\(\varphi_i\) form an equicontinuous sequence of linear maps from the
Banach space generated by \(M\) to some space \(F'_B\), where \(B\) is a
weakly closed equicontinuous disk in \(F'\). When \(F\) is also nuclear
(it even suffices for it to be a Schwartz space), we can also take
\((x'_i)\) and \((y'_i)\) in the above to be fixed, and
\(\lambda=(\lambda_i)\) to be \(\lambda=\varphi(u)\), where \(\varphi\)
is a continuous linear map from the Banach space generated by \(M\) to
the space \(\ell^\varepsilon\); we can even assume that the restriction
of \(u\) to \(M\) endowed with the simple-convergence topology to be
continuous, which gives another useful representation of equicontinuous
bilinear forms on \(E\times F\) that tend to \(0\) (in the sense of
simple convergence). Finally, we note that, if \(M\) is a set of
endomorphisms of a nuclear space \(E\), with each one sending a fixed
neighbourhood of \(0\) to a fixed bounded subset of \(E\), then the map
that sends each \(u\in M\) to the unordered sequence of its eigenvalues
is a continuous map from \(M\), endowed with the simple-convergence
topology, to the space of unordered rapidly decreasing sequences
(endowed with its natural metrisable topology): this claim can be made
explicit in the same way as the analogous claim found after
\hyperref[section-2-theorem-8]{Theorem 8}.

\subsection{\texorpdfstring{On the spaces \(E\hat{\otimes}F\) with \(E\)
of type \((\mathscr{DF})\) and \(F\) of type
\((\mathscr{F})\)}{On the spaces E\textbackslash hat\{\textbackslash otimes\}F with E of type (\textbackslash mathscr\{DF\}) and F of type (\textbackslash mathscr\{F\})}}\label{section-2-appendix}

{[}PTT, chap.~2, §4{]}

The theory of duality for spaces of the form \(E\hat{\otimes}F\), which
is very simple when \(E\) is nuclear and \(E\) and \(F\) are both of
type \((\mathscr{F})\), or both of type \((\mathscr{DF})\) (see
\hyperref[section-2-theorem-6]{Theorem 6}), becomes more complicated in
the general case. We can show, when \(E\) is nuclear, that the strong
dual \(\mathrm{B}(E,F)\) of \(E\hat{\otimes}F\) can be identified with a
dense vector subspace of \(E'\bar{\otimes}F'\), but with a topology that
is a priori \emph{less fine} (and sometimes strictly less fine) than the
topology induced by \(E'\bar{\otimes}F'\). When \(F\) is reflexive (for
example), these topologies coincide; and when, further, \(E\) and \(F\)
are both of type \((\mathscr{F})\), or both of type \((\mathscr{DF})\),
then \(\mathrm{B}(E,F)\) is a \((\mathscr{LF})\) space, in the general
sense (the generalised inductive limit of a sequence of spaces of type
\((\mathscr{F})\), cf.~{[}PTT, Introduction, IV{]}). But often, even if
\(E\) and \(F\) are both nuclear, with \(E\) of type \((\mathscr{DF})\)
and \(F\) of type \((\mathscr{F})\), the dual \(\mathrm{B}(E,F)\) of
\(E\hat{\otimes}F\) is distinct from \(E'\bar{\otimes}F'\),
i.e.~\(\mathrm{B}(E,F)\) is not complete (a fortiori \(E\hat{\otimes}F\)
is not then bornological). \oldpage{111} For example, if \(E\) is the
direct topological sum of a sequence of lines, then the disagreeable
situation described above arises whenever \(F\) is a non-normable
\((\mathscr{F})\) space that admits a continuous true norm. In {[}PTT,
chap.~2, §4, no. 1, prop. 14{]}, we study a more general case, implying,
for example, that, for the space
\(\mathscr{E}'\hat{\otimes}\mathscr{E}=\mathrm{L}(\mathscr{E},\mathscr{E})\)
(where \(\mathscr{E}=\mathscr{E}(\mathbb{R}^n)\), which is a Schwartz
space), the same troubles arise. Other important analogous spaces, such
as
\(\mathscr{S}'\hat{\otimes}\mathscr{S}=\mathrm{L}(\mathscr{S},\mathscr{S})\)
(where \(\mathscr{S}=\mathscr{S}(\mathbb{R}^n)\) is the space of
``infinitely differentiable rapidly decreasing'' functions on
\(\mathbb{R}^n\)) are, however, bornological; but we do not have simple
criteria that allow us to recognise, in general, if \(E\hat{\otimes}F\)
is bornological or not. Without supposing \(E\) or \(F\) to be nuclear,
if \(E\) is of type \((\mathscr{DF})\) and \(F\) of type
\((\mathscr{F})\), then we can, under hypotheses that are more than
sufficient in practice, we can confirm the equivalence of the following
conditions:

\begin{enumerate}
\def\labelenumi{\alph{enumi}.}
\tightlist
\item
  \(E\hat{\otimes}F\) is bornological;
\item
  \(E\hat{\otimes}F\) is barrelled;
\item
  \(\mathrm{B}(E,F)\) is complete;
\item
  \(\mathrm{B}(E,F)\) is quasi-complete (or just complete for
  sequences).
\end{enumerate}

It turns out that the hypotheses necessary for this theorem are
satisfied when \(F\) is the the ``echelon'' space constructed by the
procedure of G. Köthe \autocite{7} from an increasing sequence of
positive sequences \((b^n)=(b_i^{(n)})\). When, further, \(E\) is the
dual of a nuclear echelon space, or, more generally, if \(E\) is the
inductive limit (in the general sense) of a sequence of normed spaces
\(a^{(n)}\ell^1\) (where, for a given sequence \(a^{(n)}=(a_i^{(n)})\),
we write \(a^{(n)}\ell^1\) to denote the set of sequences given by a
product of \(a^{(n)}\) with some summable sequence), then we can give an
explicit exact criterion for \(E\hat{\otimes}F\) to be bornological:
\emph{for \(E\hat{\otimes}F\) to be bornological, it is necessary and
sufficient that, for each integer \(n_0>0\), there exist an integer
\(n>0\) such that, for every integer \(m>0\) and every positive sequence
\(\lambda=(\lambda_i)\in E'\), there exists some \(R>0\) such that, for
every pair \((i,j)\) of indices, one of the two following inequalities
holds:} \[
  \lambda_i a_i^{(m)}b_j^{(m)} \leqslant Rb_j^{(n)}
  \quad\text{or}\quad
  Ra_i^{(n)}b^{(n)} \geqslant a_i^{(m)}b_j^{(n_0)}.
\] This statement is perfectly manageable in all particular cases,
despite its barbaric aspect. In {[}PTT, chap.~2, §4, no .4{]}, I apply
this criterion to an interesting class of echelon spaces, including many
interesting spaces in analysis. We find, for example, that, if \(F\) is
the space of holomorphic functions on an open subset \(U\) of the
complex plane, then \(F'\hat{\otimes}F=\mathrm{L}(F,F)\) is bornological
if \(U=\mathbb{C}\), but not bornological if \(U\) is the unit circle!

The most important consequence of this criterion is that the space
\((s')\hat{\otimes}(s)\), where \((s)\) is the space of rapidly
decreasing sequences, is bornological. \oldpage{112} Since \((s)\) is
isomorphic, by the Hermite transform, to the space
\(\mathscr{S}(\mathbb{R}^n)\) of ``infinitely differentiable rapidly
decreasing'' functions on \(\mathbb{R}^n\), the space
\(\mathscr{S}'\otimes\mathscr{S}\) is also bornological. We can, for
example, thus conclude that the spaces \((\mathscr{O}_M)\) and
\((\mathscr{O}'_C)\) of L. Schwartz are bornological {[}PTT, chap.~2,
§4, no .4{]}. So, combining this with the results of
\hyperref[section-2.3]{§2.3}, we obtain the following:

\protect\phantomsection\label{section-2-theorem-11}
\begin{itenv}{Theorem 11}
The spaces \((\mathscr{O}_M)\) and \((\mathscr{O}'_C)\) are
bornological, complete, and nuclear, and their duals are bornological
complete nuclear \((\mathscr{LF})\) spaces.

\end{itenv}

\nocite{*}
\printbibliography[heading=bibintoc,title=Bibliography]

\end{document}
